\chapter{Transformer：ミニGPT を組み立てる}\label{ch:transformer}
\bookmap{3}

\begin{chapterlead}
\textbf{ここまで}: 各トークンが文脈の中から必要な情報を自分で選んで集めてくる仕組み、注意機構を作りました（\chapref{ch:attention}）。\\[1mm]
\textbf{この章の問い}: 注意機構を部品として、文章を生成できるモデル全体をどう組み立てればよいでしょうか。
注意のほかに、何を足せばよいのでしょうか。\\[1mm]
\textbf{この章を終えると}:
\begin{itemize}
\item Transformer（デコーダのみ）の全体図を描き、各部品の役割とテンソルの形を説明できます。
\item 位置エンコーディング、MLP、残差接続、レイヤー正規化が「なぜ必要か」を、実験の結果とともに説明できます。
\item 共有ライブラリの \code{llmbook.model.GPT} を 1 行ずつ読み、パラメータ数 2,583,936 を手で数えられます。
\item 文字単位のミニGPT を CPU で 10 分足らず学習させ、損失が下がって文章らしくなっていく様子を確かめられます。
\item 貪欲法・温度・top-$k$・top-$p$ という生成の設定を使い分けられます。
\end{itemize}
\end{chapterlead}

\begin{readingguide}
\item[\rgpre] \chapref{ch:attention}の注意の式と、\chapref{ch:lm}の交差エントロピーと温度。\chapref{ch:training}の残差接続とレイヤー正規化は、\secref{sec:tf-residual}で説明し直します。
\item[\rgmust] \secref{sec:tf-overview}〜\secref{sec:tf-gpt}：全体図と三つの部品、それらを組んだ \code{GPT} クラス。本書の後半は、このクラスを使い続けます。
\item[\rgmust] \secref{sec:tf-train}と\secref{sec:tf-generate}：学習させ、文章を書かせます。本書の最初の到達点です。ぜひ手元で動かしてください。
\item[\rgskim] \secref{sec:tf-params}（パラメータを数える）と、正弦波・GELU・\secref{sec:tf-residual}の実験・\code{generate} の読解の各小節：見出しの下に、押さえる結論を書きました。
\item[\rglater] \secref{sec:tf-family}：エンコーダや BERT との違い。その言葉に出会ったときに戻れば足ります。
\end{readingguide}

前の章の実験（\secref{sec:attn-experiment}）で、注意機構を 1 層だけ持つ小さなモデルが、1〜2 桁の足し算を 96\% 解けるようになりました。
では、この部品を何層も重ねれば、漱石の文章も書けるようになるのでしょうか。
実は、注意を重ねるだけではうまくいきません。
理由は三つあります。

一つ目に、注意の計算そのものは語順を知りません（\secref{sec:attn-properties}）。
「猫が犬を追う」と「犬が猫を追う」を区別できないのでは、言語モデルになりません。
二つ目に、注意がするのは、文脈から情報を\textbf{集める}ことだけです。
集めた情報を組み合わせて「考える」部品が、別に要ります。
三つ目に、部品を何十層も重ねると、学習が不安定になります。
勾配が消えたり、値が大きくなりすぎたりするからです（\chapref{ch:training}）。

この三つの困りごとに、一つずつ答える部品があります（\tabref{tab:tf-three}）。
これらを注意と組み合わせたものが、\term{Transformer}{Transformer}\cite{vaswani2017attention}です。
本章では、まず全体図を見てから部品を一つずつ説明し、最後に本書の基盤モデル \file{base.pt} と同じ形のミニGPT を組み上げます。
そして、漱石のコーパスで学習させ、文章を生成させます。

行き先を先に見ておきましょう。
この章で組み上げるモデルに「吾輩は」の続きを書かせた、学習の前と、手元の CPU で 8 分ほど学習させた後の出力です（\secref{sec:tf-train}）。
\begin{quote}
学習前：吾輩は堅幇博耄奨蜷諷描岫賤色功室鞄伐舎盤帽搏雨媒憫憫翹笊漸卍…\\
8 分後：吾輩はあれった自分が、僕の女の室に、その晩は無論なお延の不思議に…
\end{quote}
意味はまだ通りませんが、でたらめな漢字の列が、小説らしい文字の並びに変わっています。
この変化を、\chapref{ch:nn}から作ってきた部品だけで、自分の手で起こすのがこの章の到達点です。

\begin{table}[tbp]
\centering
\caption{注意を重ねるだけでは足りない三つの理由と、それに答える部品。本章の前半は、この表を上から順にたどる。}
\label{tab:tf-three}
\small
\begin{tabular}{llll}
\toprule
注意の困りごと & 答える部品 & ひとことで言うと & 節 \\
\midrule
語順を知らない & 位置エンコーディング & 「何番目か」を入力に足す & \ref{sec:tf-position} \\
情報を集めるだけ & MLP & 集めた情報を各位置で加工する & \ref{sec:tf-mlp} \\
深く重ねると不安定 & 残差接続、レイヤー正規化 & 足し込む。大きさをそろえる & \ref{sec:tf-residual} \\
\bottomrule
\end{tabular}
\end{table}

\section{全体図：ミニGPT を上から眺める}\label{sec:tf-overview}

部品の説明に入る前に、完成品の姿を見ておきましょう。
地図を持たずに部品を一つずつ見ていくと、いま何のためにそれを作っているのかが分からなくなるからです。
この節の問いは、「トークン列を入れてから次のトークンの確率が出てくるまでに、データはどんな部品を、どんな形で通るのか」です。

\subsection*{入力から出力まで}

本書のミニGPT は、\term{デコーダのみ}{decoder-only}の Transformer です。
「デコーダ」は、もとは翻訳のモデルの「訳文を書く側」の名前でした。
ここでは「過去のトークンだけを見て、次のトークンを書く部品」という意味だと思っておけば十分です（由来は\secref{sec:tf-family}で説明します）。
\figref{fig:tf-overview}が全体図です。
前の章で足し算を学習させたモデル（\figref{fig:attn-exp-model}）は、この図のブロックが 1 個だけで、LayerNorm のないものでした。
下から上へ、次の順にデータが流れます。

\begin{enumerate}
\item \textbf{入力}: トークン ID の列 $x_1,\dots,x_T$。たとえば「\texttt{\textbackslash n12+35=}」を BPE でトークン化したものです。
\item \textbf{埋め込み}: 各トークン ID を $d$ 次元のベクトルに変えます（\chapref{ch:embedding}）。そこに「何番目か」を表す\textbf{位置のベクトル}を足します。
\item \textbf{Transformer ブロック × $L$}: 同じ形のブロックを $L$ 個重ねます。各ブロックの中には、注意（\chapref{ch:attention}）と MLP があり、どちらも残差接続でつながっています。
\item \textbf{最後の正規化}: ベクトルの大きさをそろえます。
\item \textbf{出力ヘッド}: 各位置の $d$ 次元のベクトルを、語彙の数 $V$ 個のロジットに変えます。softmax をとれば、次のトークンの確率分布になります。
\end{enumerate}

3 の\textbf{残差接続}は、部品の出力を入力に足し込むつなぎ方です。
4 の正規化は\textbf{レイヤー正規化}で、ベクトルの成分の平均と大きさをそろえ直します（図とコードでは LayerNorm と書きます）。
どちらも\chapref{ch:training}の部品ですが、\secref{sec:tf-residual}で式からもう一度説明します。
\chapref{ch:training}をまだ読んでいない人は、ここでは名前だけ見ておいてください。

\begin{figure}[tbp]
\centering
\begin{tikzpicture}[x=1mm,y=0.95mm,
  w/.style={font=\footnotesize,minimum width=24mm,minimum height=5.4mm,inner ysep=0.4mm},
  wide/.style={w,minimum width=50mm},
  io/.style={tok,font=\footnotesize,minimum height=5.4mm},
  num/.style={draw=bookgray,circle,inner sep=0.8pt,font=\scriptsize\sffamily,text=bookgray,fill=white}]
  \draw[bookblue,thick,rounded corners=3pt,fill=figblue!25] (-38,22) rectangle (40,74.5);
  \node[num] at (-34.6,70.9) {3};
  \node[lbl,anchor=north west,align=left] at (-32.8,73.4) {Transformer\\ブロック $\times L$};
  \node[io] (in) at (-20,0) {トークン ID $x_1,\dots,x_T$\ \ $(B,T)$};
  \node[io,fill=figgray] (ps) at (21,0) {位置の番号 $0,\dots,T{-}1$};
  \node[box,w] (te) at (-20,8.5) {トークン埋め込み};
  \node[orangebox,w] (pe) at (21,8.5) {位置埋め込み};
  \node[op] (p0) at (0,16) {$+$};
  \node[graybox,w] (ln1) at (0,30) {LayerNorm};
  \node[redbox,w] (att) at (0,38) {因果自己注意};
  \node[op] (p1) at (0,45.6) {$+$};
  \node[graybox,w] (ln2) at (0,54.8) {LayerNorm};
  \node[greenbox,w] (mlp) at (0,62.8) {MLP};
  \node[op] (p2) at (0,70.2) {$+$};
  \node[graybox,wide] (lnf) at (0,82.5) {最後の LayerNorm};
  \node[purplebox,wide] (head) at (0,90.5) {出力ヘッド（$d\to V$）};
  \node[io,minimum width=50mm] (out) at (0,99) {ロジット $\vz_1,\dots,\vz_T$\quad $(B,T,V)$};
  \draw[arr] (in) -- (te); \draw[arr] (ps) -- (pe);
  \draw[arr] (te.east) -- (p0); \draw[arr] (pe.west) -- (p0);
  \draw[arr] (p0) -- (ln1); \draw[arr] (ln1) -- (att); \draw[arr] (att) -- (p1);
  \draw[arr] (p1) -- (ln2); \draw[arr] (ln2) -- (mlp); \draw[arr] (mlp) -- (p2);
  \draw[arr] (p2) -- (lnf); \draw[arr] (lnf) -- (head); \draw[arr] (head) -- (out);
  \draw[arr,bookblue,very thick] (0,24.8) -- (19,24.8) |- (p1);
  \draw[arr,bookblue,very thick] (0,50) -- (19,50) |- (p2);
  \fill[bookblue] (0,24.8) circle (0.6mm); \fill[bookblue] (0,50) circle (0.6mm);
  \node[lbl,anchor=west] at (20,36) {残差接続};
  \node[lbl,anchor=west] at (20,61) {残差接続};
  \node[lbl,anchor=east,align=right] at (-14,34) {通信\\（位置の間）};
  \node[lbl,anchor=east,align=right] at (-14,58.8) {計算\\（各位置で）};
  \node[lbl,anchor=west] at (2,19.7) {$(B,T,d)$};
  \node[lbl,anchor=west] at (2,76.8) {$(B,T,d)$};
  \node[num] at ($(in.west)+(-3,0)$) {1};
  \node[num] at ($(te.west)+(-3,0)$) {2};
  \node[num] at ($(lnf.west)+(-3,0)$) {4};
  \node[num] at ($(head.west)+(-3,0)$) {5};
\end{tikzpicture}
\caption{デコーダのみの Transformer（ミニGPT）の全体図。下から上へデータが流れ、丸の数字は本文の手順 1〜5 に対応する。ブロックの中の右側の線が残差接続で、注意と MLP の出力は「元のベクトルに足し込まれる」。
本書の基盤モデルでは $L=4$、$d=192$、$V=4096$、$T\le 128$。}
\label{fig:tf-overview}
\end{figure}

図で見てほしいのは、次の 2 点です。
第一に、埋め込みを足した直後から最後の正規化の手前まで、テンソルの形は $(B,T,d)$ のまま変わりません。
形が変わるのは、入口（$(B,T)\to(B,T,d)$）と出口（$(B,T,d)\to(B,T,V)$）だけです。
入口では、1 個の整数だったトークン ID が $d$ 次元のベクトルに変わるので、軸が 1 本増えます。
出口では、$d$ 次元のベクトルが語彙 $V$ 個分の点数に変わるので、最後の軸の長さが $V$ になります。
だからこそ、同じ形のブロックを何個でも重ねられるのです。
第二に、ブロックの中の右側の線（残差接続）です。
注意と MLP は、入ってきたベクトルを\textbf{置き換える}のではなく、\textbf{足し込む}ものを計算します。
この「足し込む」という設計が、深く重ねても学習できる鍵になります（\secref{sec:tf-residual}）。

\subsection*{形の追跡表}

各段階のテンソルの形を\tabref{tab:tf-shapes}にまとめました。
$B$ はバッチサイズ、$T$ は系列長、$d$ は埋め込み次元、$V$ は語彙の大きさです。
コードでは $d$ を \code{C}（\code{n_embd}）と書くことも思い出しておきましょう（\chapref{ch:attention}）。

\begin{table}[tbp]
\centering
\caption{ミニGPT の各段階のテンソルの形。右端は基盤モデル（$d=192$, $V=4096$, $H=4$）で $B=2$, $T=8$ のときの値。形が変わるのは入口と出口だけである。}
\label{tab:tf-shapes}
\small
\begin{tabular}{llll}
\toprule
段階 & コード & 形 & 例 \\
\midrule
入力（トークン ID） & \code{idx} & $(B,T)$ & $(2,8)$ \\
トークン埋め込み & \code{tok_emb(idx)} & $(B,T,d)$ & $(2,8,192)$ \\
位置埋め込み & \code{pos_emb(pos)} & $(T,d)$ & $(8,192)$ \\
足したもの（ブロックへの入力） & \code{x} & $(B,T,d)$ & $(2,8,192)$ \\
注意の中のヘッドごとの $\mQ,\mK,\mV$ & \code{q, k, v} & $(B,H,T,\dhead)$ & $(2,4,8,48)$ \\
注意の重み & \code{att} & $(B,H,T,T)$ & $(2,4,8,8)$ \\
MLP の中間 & \code{fc(x)} & $(B,T,4d)$ & $(2,8,768)$ \\
ブロックの出力 & \code{block(x)} & $(B,T,d)$ & $(2,8,192)$ \\
ロジット & \code{logits} & $(B,T,V)$ & $(2,8,4096)$ \\
\bottomrule
\end{tabular}
\end{table}

位置埋め込みだけ形が $(T,d)$ で、バッチの軸がないことに気づいたでしょうか。
位置のベクトルは「何番目か」だけで決まり、どの文でも同じだからです。
足すときには、ブロードキャスト（\chapref{ch:pytorch}）でバッチの全員に同じものが足されます。

表の中ほどの 3 行は、ブロックの\textbf{中}の形です。
注意の中では $d$ 次元を $H$ 個のヘッドに等分するので $\dhead=d/H=192/4=48$、MLP の中では一時的に $4d=768$ になり、出るときには $(B,T,d)$ に戻ります。

\begin{quiz}
基盤モデル（$L=4$, $d=192$, $V=4096$）に、$B=4$ 個の文をそれぞれ $T=16$ トークンずつ入れます。
(1) 3 個目のブロックを出た直後、(2) MLP の中間、(3) ロジット、のテンソルの形を、\tabref{tab:tf-shapes}を見ながら答えてください。
\end{quiz}

\paragraph{答え}
(1) は $(4,16,192)$ です。ブロックを何個通っても形は $(B,T,d)$ のままです。
(2) は $(4,16,768)$ で、最後の軸だけが $4d$ に広がります。
(3) は $(4,16,4096)$ です。$16$ 個の位置のそれぞれに、語彙 $4096$ 個分のロジットが出ます。

この節で分かったのは、ミニGPT が「入口で形を $(B,T,d)$ にし、同じ形のブロックを重ね、出口で $(B,T,V)$ にする」という単純な構造だということです。
次の節から、図の下から順に部品を見ていきます。

\section{位置エンコーディング：語順を教える}\label{sec:tf-position}

全体図のいちばん下には、トークン埋め込みと並んで「位置埋め込み」がありました。
前の章で、注意は語順を知らないことを確かめました（\secref{sec:attn-properties}）。
この節の問いは、「語順の情報を、どうやってモデルに渡せばよいか」です。

\subsection*{何に困っているか}

注意の出力は、キーとバリューの\textbf{集合}だけで決まり、並び順には関係しません。
たとえば計算ドリルの「\code{12+35=}」と「\code{21+53=}」を考えます。
最後の「\code{=}」から見ると、どちらも同じ数字 1, 2, 3, 5 と記号 \code{+} が過去にある、というだけになってしまいます。
しかし答えは 47 と 74 で違います。
語順が分からなければ、足し算すらできないのです。

\subsection*{いちばん素直な解決：位置ごとのベクトルを足す}

解決の考え方は単純です。
「何番目のトークンか」という情報を、ベクトルにして埋め込みに足してしまえばよいのです。
位置 $t$ にあるトークン $x_t$ の入力ベクトルを、次のように作ります。
\begin{equation}
\vh_t^{(0)} = \ve_{x_t} + \vp_t \qquad (t=1,\dots,T)
\label{eq:tf-input}
\end{equation}
ここで $\ve_{x_t}\in\R^d$ はトークン $x_t$ の埋め込み（\chapref{ch:embedding}）、$\vp_t\in\R^d$ は位置 $t$ のベクトルです。
右肩の $(0)$ は「ブロックに入る前（0 層目）」を表します。
\figref{fig:tf-posadd}を見てください。
同じ「\code{1}」というトークンでも、1 番目にあれば $\ve_{\code{1}}+\vp_1$、4 番目にあれば $\ve_{\code{1}}+\vp_4$ と違うベクトルになります。
クエリとキーは、この入力ベクトルに行列を掛けて作るので、位置の情報がクエリとキーにも混ざります。
これで、「すぐ前のトークン」「3 つ前のトークン」を区別して見られるようになります。
このように、位置の情報をベクトルで表してモデルに渡す仕組みを\term{位置エンコーディング}{positional encoding}と呼びます。

\begin{figure}[tbp]
\centering
\begin{tikzpicture}[x=14mm,y=9.5mm,
  v/.style={minimum width=11.5mm,minimum height=6mm,font=\scriptsize,inner sep=1pt}]
  \foreach \i/\t/\s in {1/1/orangebox,2/2/box,3/+/box,4/1/orangebox,5/5/box,6/=/box} {
    \node[tok,minimum width=8mm] (x\i) at (\i,3) {\texttt{\t}};
    \node[\s,v] (e\i) at (\i,2) {$\ve_{\texttt{\t}}$};
    \node[graybox,v] (p\i) at (\i,1) {$\vp_{\i}$};
    \node[greenbox,v] (h\i) at (\i,0) {$\ve_{\texttt{\t}}{+}\vp_{\i}$};
    \draw[arr] (x\i) -- (e\i);
    \node[font=\scriptsize] at (\i,1.5) {$+$};
    \draw[arr] (p\i) -- (h\i);
  }
  \node[lbl,anchor=east] at (0.45,3) {トークン};
  \node[lbl,anchor=east] at (0.45,2) {トークン埋め込み};
  \node[lbl,anchor=east] at (0.45,1) {位置埋め込み};
  \node[lbl,anchor=east] at (0.45,0) {入力 $\vh^{(0)}_t$};
\end{tikzpicture}
\caption{位置埋め込みを足す（式\eqref{eq:tf-input}）。説明のため 1 文字を 1 トークンとした。
2 か所の「\texttt{1}」は、トークン埋め込み（橙）は同じだが、足される位置のベクトルが $\vp_1$ と $\vp_4$ で違うので、ブロックに入るベクトル（緑）は別のものになる。}
\label{fig:tf-posadd}
\end{figure}

\begin{definition}[位置エンコーディング（学習型の位置埋め込み）]
文脈長の上限を $T_{\max}$ として、行列 $\mP\in\R^{T_{\max}\times d}$ を用意し、その $t$ 行目 $\vp_t$ を位置 $t$ のベクトルとする。
入力は式\eqref{eq:tf-input}のようにトークン埋め込みに $\vp_t$ を足して作る。
$\mP$ はほかのパラメータと同じく、勾配降下法で学習する。
これを\term{位置埋め込み}{position embedding}と呼ぶ。
\end{definition}

トークン埋め込みが「トークン ID で行を引く表」だったのと同じく、位置埋め込みは「位置の番号で行を引く表」です。
コードでも、どちらも \code{nn.Embedding} です（\secref{sec:tf-gpt}）。
基盤モデルでは $T_{\max}=128$（\code{block_size}）、$d=192$ なので、$\mP$ は $128\times 192$ の行列です。
表には 128 行しかないので、129 番目の位置には、引く行がありません。

\begin{caution}{足したら情報が混ざってしまわないか}
トークンのベクトルと位置のベクトルを「足す」と、両者が区別できなくなりそうに思えます。
しかし $d=192$ 次元もあれば、学習によって「トークンの情報を主に載せる方向」と「位置の情報を主に載せる方向」を使い分けられます。
足し算は、2 つのベクトルを縦に並べて（連結して）から線形変換するのと、表現力の点でほぼ変わりません。
連結したベクトルに行列 $[\mA\ \mB]$ をかけると
$[\mA\ \mB]\begin{bmatrix}\ve\\ \vp\end{bmatrix}=\mA\ve+\mB\vp$
となり、結局は「それぞれを変換してから足したもの」になるからです。
しかも、トークン埋め込みも位置埋め込みも学習される表なので、$\mA\ve$ にあたるものを $\ve$ として、$\mB\vp$ にあたるものを $\vp$ として、初めから覚えてしまえます。
足し算のほうが次元が増えないので、こちらが標準になっています。
\end{caution}

\subsection*{確かめる：位置埋め込みがあると語順が分かる}

本当に語順が区別できるようになったかを、小さなモデルで確かめます。
トークン列「5 7 9」と、先頭の 2 つを入れ替えた「7 5 9」を入れ、最後の位置のロジットを比べます。

\snippet{位置埋め込みの有無と、語順の入れ替えへの反応}{code:tf-pos-check}{ch10/pos_encoding.py}{pos-check}

1〜4 行目で、1 層・16 次元の小さなモデル（学習前）と、比べる 2 つの入力を作ります（5, 7, 9 は適当に選んだトークン ID です）。
\code{GPT} は\secref{sec:tf-gpt}で読むクラスで、ここでは「トークン ID を入れると（ロジット, 損失）の組を返す箱」として使います。
2 行目の \code{.eval()}（評価用のモード）と 5 行目の \code{torch.no_grad()}（勾配の記録をとらない。\secref{sec:pt-autograd}）は、学習をせず出力を見るだけのときの指定です。
6〜7 行目の \code{model(a)[0]} がそのロジットで、続く \code{[0, -1]} で最後の位置を取り出し、差の絶対値の最大を表示します。
8 行目で位置埋め込みの重みを \code{zero_()} ですべて 0 にして位置の情報を消し、9〜10 行目で同じ比較をもう一度します。

\outputfile{ch10/out/pos_check.txt}

位置埋め込みがあるときは、ロジットに $3\times10^{-3}$ 程度の差が出ます（学習前のモデルなので差は小さいですが、0 ではありません）。
学習前でも差が出るのは、位置埋め込みの表が乱数で初期化されていて、行ごとに違うベクトルが入っているからです。
位置埋め込みを消すと、差は $3\times 10^{-8}$、つまり浮動小数点の丸め誤差の大きさになり、2 つの入力は区別できなくなります。
ぴったり 0 にならないのは、足す順番が変わると、浮動小数点の最後の桁がわずかにずれるからです。
最後の「9」から見ると、「5」と「7」のどちらが先だったのかが分からないのです。
位置のベクトルを足しただけで、語順が見えるようになりました。

\subsection*{原論文の正弦波の位置エンコーディング}
\skimread{押さえる結論は「原論文は位置のベクトルを sin と cos の式で作った。本書と GPT-2 は、学習する表を使う」です。式の細部は 2 周目で足ります。}

Transformer の原論文\cite{vaswani2017attention}は、位置のベクトルを学習せず、決まった式で作りました（本書のモデルでは使いませんが、よく目にするので仕組みだけ見ておきます）。
位置 $t$（ここだけ 0 始まり）の $\vp_t$ の第 $2i$ 成分と第 $2i+1$ 成分（$i=0,1,\dots,d/2-1$）を
\begin{equation}
p_{t,2i} = \sin\!\left(\frac{t}{10000^{2i/d}}\right),\qquad
p_{t,2i+1} = \cos\!\left(\frac{t}{10000^{2i/d}}\right)
\label{eq:tf-sinusoid}
\end{equation}
とします。
成分の組 $(2i, 2i+1)$ ごとに、$\sin$ と $\cos$ の波を割り当てるわけです（10000 は原論文が選んだ定数です）。
$i=0$ の組は、分母が $10000^{0}=1$ なので周期 $2\pi\approx 6.3$ で速く回り、$i$ が大きくなるほど周期が長くなります。
たとえば $d=16$、位置 $t=1$ では、$i=0$ の成分が $\sin(1)=0.8415$ です。
$i=1$ の成分は、分母が $10000^{2/16}\approx3.162$ なので、$\sin(1/3.162)=\sin(0.3162)=0.3110$ です。

\begin{figure}[tbp]
\centering
\begin{tikzpicture}
\begin{axis}[width=0.92\linewidth,height=0.31\linewidth,xlabel={位置 $t$},ylabel={値},
  xmin=0,xmax=63,ymin=-1.15,ymax=1.15,legend columns=4,
  legend style={at={(0.5,1.03)},anchor=south,font=\footnotesize}]
\addplot[bookblue,thick,no marks] table[col sep=comma,x=pos,y=d0]{code/ch10/out/sinusoid.csv};
\addlegendentry{成分 0（$i=0$）}
\addplot[bookorange,thick,no marks] table[col sep=comma,x=pos,y=d2]{code/ch10/out/sinusoid.csv};
\addlegendentry{成分 2（$i=1$）}
\addplot[bookgreen,thick,no marks] table[col sep=comma,x=pos,y=d4]{code/ch10/out/sinusoid.csv};
\addlegendentry{成分 4（$i=2$）}
\addplot[bookpurple,thick,no marks] table[col sep=comma,x=pos,y=d6]{code/ch10/out/sinusoid.csv};
\addlegendentry{成分 6（$i=3$）}
\end{axis}
\end{tikzpicture}
\caption{正弦波の位置エンコーディング（式\eqref{eq:tf-sinusoid}、$d=16$）の偶数番目の成分（整数の位置での値を線で結んだ）。成分の番号が大きいほど波がゆっくり変わる。
時計の秒針・分針・時針のように、速い波と遅い波の組み合わせで、位置ごとに異なるベクトルができる。}
\label{fig:tf-sinusoid}
\end{figure}

\figref{fig:tf-sinusoid}を見てください。
成分 0（青）は数位置ごとに上下しますが、成分 6（紫）は 64 位置かけてもまだ半周しません。
時計の秒針・分針・時針の組で時刻が一つに決まるのと同じく、速い波と遅い波の組で位置が一つに決まります。
正弦波の利点は、学習するパラメータがいらないことと、学習中に見たより長い位置にも式で値を作れることです。
一方、学習型は、データに合わせて位置の表し方を自由に決められます。
GPT-2\cite{radford2019language} は学習型を選び、本書もそれに従います。
現在の LLM の多くは、どちらでもない\textbf{回転}による位置表現（RoPE）を使います（\chapref{ch:modern-arch}）。

この節で分かったのは、語順の情報は「位置ごとのベクトルを埋め込みに足す」だけで渡せるということです。

\section{MLP：各位置で「考える」部品}\label{sec:tf-mlp}

位置の情報が加わり、注意は「どこに何があるか」を見て情報を集められるようになりました。
しかし、注意の出力はバリューの重み付き平均、つまり\textbf{集めてきた情報を混ぜたもの}にすぎません。
この節の問いは、「集めた情報を組み合わせて、新しい情報を作る部品は何か」です。

\subsection*{注意だけでは足りない理由}

計算ドリル「\code{25+48=}」の続きを書くことを考えます。
十の位の「7」を書くには、「2」と「4」を集めてくるだけでは足りません。
2 と 4 を足し、さらに一の位の 5 と 8 の和が 10 以上かどうか（繰り上がり）を判定して、1 を足す必要があります。
これは「集める」ではなく「計算する」仕事です。

注意の計算を振り返ると、非線形な部分は softmax で重みを決めるところだけでした。
バリューに対しては、重み付きの和という\textbf{線形な}操作しかしていません。
線形な操作をいくら重ねても線形のままなので（\chapref{ch:nn}）、繰り上がりの判定のような「しきい値で分ける」計算は苦手です。
そこで、各位置に、非線形な活性化関数を持つ MLP（\chapref{ch:nn}）を置きます。

\subsection*{位置ごとの MLP}

Transformer の MLP は、2 層の全結合ネットワークです。
位置 $t$ のベクトル $\vu_t\in\R^d$ に対して、
\begin{equation}
\mathrm{MLP}(\vu_t) = \mW_2^\T\,\gelu\!\left(\mW_1^\T \vu_t\right),\qquad
\mW_1\in\R^{d\times 4d},\ \mW_2\in\R^{4d\times d}
\label{eq:tf-mlp}
\end{equation}
を計算します。
転置の付いた書き方は、\chapref{ch:attention}の $\vq_s=\mW_q^\T\vx_s$ と同じ約束です（$\mW_1$ は「入力の次元 × 出力の次元」の形）。
1 層目で $d$ 次元を 4 倍の $4d$ 次元（基盤モデルでは 768）に広げ、活性化関数 GELU（すぐ後で説明します）を通し、2 層目で $d$ 次元に戻します（\figref{fig:tf-mlp-shape}）。
いったん広げるのは、活性化関数を通す「しきい値の判定」を、たくさん並べて一度に行うためです。
4 倍という値は原論文\cite{vaswani2017attention}以来の慣習で、理論的に決まった値ではありません。
\code{llmbook} ではバイアスを省いています（GPT-2 にはありますが、なくても学習でき、パラメータも数えやすくなります）。

ここで大事なのは、\textbf{同じ MLP（同じ $\mW_1,\mW_2$）を、すべての位置に別々に適用する}ことです。
MLP は位置 $t$ のベクトルだけを見て、ほかの位置のベクトルは見ません。
位置の間で情報をやりとりするのは注意だけ、各位置で計算するのは MLP だけ、という分業になっています。
同じ重みを使い回すのは、どの位置でも「集まった情報を加工する」という同じ種類の仕事をするからです（位置ごとに別の重みを持つと、パラメータが系列長に比例して増えてしまいます）。

\begin{figure}[tbp]
\centering
\begin{tikzpicture}[x=13mm,y=10mm]
  \foreach \i/\t in {1/2,2/5,3/+,4/4,5/8,6/=} {
    \node[tok,minimum width=8mm] (x\i) at (\i,0) {\texttt{\t}};
    \node[redbox,minimum width=9mm,minimum height=6mm,font=\scriptsize] (a\i) at (\i,1.5) {注意};
    \node[greenbox,minimum width=9mm,minimum height=6mm,font=\scriptsize] (m\i) at (\i,3) {MLP};
    \draw[arr] (x\i) -- (a\i);
    \draw[arr] (a\i) -- (m\i);
  }
  \node[draw=bookaccent,thick,rounded corners,fit=(a1)(a6),inner sep=2pt,label={[lbl]left:通信}] {};
  \foreach \s in {1,2,4,5} { \draw[->,bookaccent,thick] (x\s.north) to[out=18,in=200] ([xshift=-1.5mm]a6.south); }
  \node[lbl,anchor=east] at (0.42,3) {計算};
  \node[lbl,anchor=west,align=left] at (6.6,3) {同じ重みの MLP を\\位置ごとに別々に};
  \node[lbl,anchor=west,align=left] at (6.6,1.5) {「=」の位置が\\2, 5, 4, 8 を集める};
\end{tikzpicture}
\caption{注意と MLP の分業。注意（赤）は位置の間で情報をやりとりする「通信」で、赤い曲線は「=」の位置が数字を集める様子（ほかの位置も、それぞれ自分より前から集めている）。
MLP（緑）は各位置で、集まった情報から新しい情報（たとえば繰り上がりの有無）を作る「計算」。MLP は隣の位置を見ない。}
\label{fig:tf-comm-comp}
\end{figure}

\figref{fig:tf-comm-comp}は、この分業を「\code{25+48=}」の例で描いたものです。
赤い曲線が注意による情報の集約で、「=」の位置に 2, 5, 4, 8 の情報が集まります。
その上の緑の MLP は、「=」の位置に集まった情報だけを使って、答えの十の位を計算するための特徴を作ります。
Karpathy の講義\cite{karpathy2022zerotohero}では、この分業を「注意は\textbf{通信}（communication）、MLP は\textbf{計算}（computation）」と表現しています。
ブロックを重ねると「通信 → 計算 → 通信 → 計算 …」が繰り返されるので、集めた結果をもとに、さらに別の場所から情報を集めにいけます。

\subsection*{活性化関数 GELU}
\skimread{押さえる結論は「GELU は、ReLU の角を丸めた活性化関数」です。式の中身は 2 周目で足ります。ただし、小節の終わりの MLP のコード（\coderef{code:tf-mlp}）とその説明は読んでください。}

式\eqref{eq:tf-mlp}の活性化関数には、ReLU（\chapref{ch:nn}）ではなく GELU\cite{hendrycks2016gelu}を使います。
\begin{equation}
\gelu(u) = u\,\Phi(u)
\label{eq:tf-gelu}
\end{equation}
$\Phi(u)$ は、平均 0・分散 1 の正規分布に従う乱数が $u$ 以下になる確率で、$u$ が大きな正の数なら 1、大きな負の数なら 0 に近づきます。
したがって GELU は、正の大きな入力はほぼそのまま通し（$\gelu(3)=3\times0.9987\approx 3.00$）、負の大きな入力はほぼ 0 にします（$\gelu(-3)=-3\times0.0013\approx -0.004$）。
ここまでは ReLU と同じです。
違うのは 0 の近くで、たとえば $\gelu(-0.5)=-0.5\times 0.3085\approx-0.154$ のように、少しだけ負の値を通し、角のない滑らかな曲線になります（\figref{fig:tf-gelu}。数値は \file{ch10/gelu_curve.py} で計算しました）。
ReLU との差は小さく、「GPT-2 が使ったので本書も使う」と考えて構いません。
現在の LLM では、GELU の代わりに SwiGLU という MLP の変種がよく使われます（\chapref{ch:modern-arch}）。

\begin{figure}[tbp]
\begin{minipage}[b]{0.52\linewidth}
\centering
\begin{tikzpicture}[x=1mm,y=1mm,
  bar/.style={minimum height=4.5mm,font=\scriptsize,inner sep=1pt}]
  \node[box,bar,minimum width=14mm] (a) at (0,0) {$d=192$};
  \node[greenbox,bar,minimum width=56mm] (b) at (0,11) {$4d=768$ 次元に広げる};
  \node[greenbox,bar,minimum width=56mm] (c) at (0,22) {負の成分をほぼ 0 にする};
  \node[box,bar,minimum width=14mm] (d) at (0,33) {$d=192$};
  \draw[arr] (a) -- node[lbl,right] {\code{fc}（$\mW_1$）} (b);
  \draw[arr] (b) -- node[lbl,right] {GELU} (c);
  \draw[arr] (c) -- node[lbl,right] {\code{proj}（$\mW_2$）} (d);
  \node[lbl,anchor=east] at (-9,0) {入力 $\vu_t$};
  \node[lbl,anchor=east] at (-9,33) {出力};
\end{tikzpicture}
\caption{MLP の中の形（基盤モデル）。下から上へ進む。いったん 4 倍に広げ、GELU を通してから元の幅に戻す。}
\label{fig:tf-mlp-shape}
\end{minipage}\hfill
\begin{minipage}[b]{0.44\linewidth}
\centering
\begin{tikzpicture}
\begin{axis}[width=1.08\linewidth,height=0.86\linewidth,xlabel={入力 $u$},
  xmin=-4,xmax=4,ymin=-0.6,ymax=4,
  legend pos=north west,legend style={font=\scriptsize}]
\addplot[bookblue,very thick,no marks] table[col sep=comma,x=u,y=gelu]{code/ch10/out/gelu.csv};
\addlegendentry{GELU}
\addplot[bookorange,thick,dashed,no marks] table[col sep=comma,x=u,y=relu]{code/ch10/out/gelu.csv};
\addlegendentry{ReLU}
\end{axis}
\end{tikzpicture}
\caption{GELU と ReLU。違いは 0 の近くだけで、GELU は角がなく、少しだけ負の値を通す。}
\label{fig:tf-gelu}
\end{minipage}
\end{figure}

\code{llmbook} の MLP のコードは次のとおりです。

\snippet{共有ライブラリの MLP（位置ごとの 2 層ネットワーク）}{code:tf-mlp}{llmbook/model.py}{mlp}

6 行目の \code{fc} が $\mW_1$（$d\to 4d$）、7 行目の \code{proj} が $\mW_2$（$4d\to d$）で、どちらも \code{bias=False} です。
\code{nn.Linear} は入力の最後の軸にだけ作用するので（\chapref{ch:nn}）、$(B,T,d)$ のテンソルを入れれば、$B\times T$ 個の位置それぞれに同じ重みが適用されます。
「位置ごとに別々に」は、この性質だけで実現されているのです。
なお、\code{nn.Linear} は重みを「出力 × 入力」の形で保存するので、\code{fc.weight} の形 $(4d,d)$ は式の $\mW_1^\T$ に当たります。
11 行目が式\eqref{eq:tf-mlp}そのもので、\code{fc} → \code{F.gelu} → \code{proj} の順に通し、最後にドロップアウト（\chapref{ch:training}。基盤モデルでは率 0 なので何もしない）をかけます。
\code{proj} という名前は注意の出力射影と同じで、どちらも「残差に足し込む直前の層」です。この名前は初期化（\secref{sec:tf-gpt}）で使われます。

この節で分かったのは、注意が集めた情報を、各位置の MLP が非線形に組み合わせる、という分業です。

\section{深く積むための工夫：残差接続とレイヤー正規化}\label{sec:tf-residual}

注意（通信）と MLP（計算）がそろいました。
この 2 つを 1 組にして何層も重ねれば、「集めて考える」を何度も繰り返せます。
しかし、層を深く重ねたネットワークは、そのままでは学習がうまく進みません。
\secref{sec:nn-deep}の実験では、ReLU の MLP を隠れ層 10 枚にしただけで、損失が何も学んでいないときの値 $\log 3\approx1.099$ から動かなくなりました。
逆伝播では層の数だけ掛け算が重なるので、勾配が途中で消えたり、爆発したりするのでした。
\chapref{ch:training}は、これを初期化・レイヤー正規化・残差接続で解き、50 枚の MLP を 5 回とも学習させました。

この節の問いは、「注意と MLP を何層も重ねても学習が進むようにするには、どうつなげばよいか」です。
答えは、\chapref{ch:training}の二つの道具、残差接続とレイヤー正規化を、決まった順序で組み込むことです。
\chapref{ch:training}を後回しにした人も読めるように、式と理由をここでもう一度説明します。

\subsection*{残差接続：置き換えずに足し込む}

\term{残差接続}{residual connection}\cite{he2016deep}は、ある部品 $F$ の出力で入力を置き換えるのではなく、入力に足し込むつなぎ方でした（\secref{sec:training-residual}）。
\begin{equation}
\vh \leftarrow \vh + F(\vh)
\label{eq:tf-residual}
\end{equation}
Transformer ブロックでは、注意と MLP のそれぞれに残差接続を付けます。
なぜこれで深くしても学習できるのかを、勾配でおさらいします。
式\eqref{eq:tf-residual}の右辺では、$\vh$ が「そのまま足される道」と「$F$ の入力」の 2 か所で使われています。
逆伝播では、出口側から届いた勾配 $\vg$ が足し算のところで 2 つの道に分かれるので、$\vh$ に届く勾配は
\begin{equation}
\underbrace{\vg}_{\text{近道を通る分}}+\underbrace{(F\text{ を逆向きに通った }\vg)}_{F\text{ を通る分}}
\label{eq:tf-residual-grad}
\end{equation}
です（式\eqref{eq:training-residual-grad}と同じです）。
2 か所で使われた値の勾配は足す、という\chapref{ch:nn}の規則です（式\eqref{eq:nn-fanout}）。
足し算は勾配をそのまま渡すので、第 1 項は何も掛けられない $\vg$ のままです。
足し込みを何回重ねても、この「近道」は出口から入口まで 1 本につながっています。
したがって、$F$ を通る道で勾配がどれだけ小さくなっても、損失の勾配は近道を通って、入口近くの層まで届きます。
残差接続がなければ、勾配はすべての $F$ を順に通るしかなく、層が深いほど消えたり爆発したりしやすくなります。

\subsection*{残差ストリームという見方}

残差接続のおかげで、ブロックの中のデータの流れを、別の見方で眺められます。
埋め込みを足した直後のベクトルを $\vh^{(0)}$ とし、$\ell$ 番目のブロックの注意と MLP が足し込む量をそれぞれ $\va^{(\ell)}$、$\vm^{(\ell)}$ とすると、最後のブロックの出力は
\begin{equation}
\vh^{(L)} = \vh^{(0)} + \sum_{\ell=1}^{L}\left(\va^{(\ell)} + \vm^{(\ell)}\right)
\label{eq:tf-stream}
\end{equation}
と書けます（式\eqref{eq:tf-residual}を $2L$ 回くり返して、足し込んだ量を並べただけです）。
1 番目のブロックを出ると $\vh^{(1)}=\vh^{(0)}+\va^{(1)}+\vm^{(1)}$、2 番目を出ると $\vh^{(2)}=\vh^{(1)}+\va^{(2)}+\vm^{(2)}=\vh^{(0)}+\va^{(1)}+\vm^{(1)}+\va^{(2)}+\vm^{(2)}$ で、これを $L$ 番目まで続けます。
つまり、各位置には 1 本の $d$ 次元の「流れ」があり、注意と MLP はそこから読み出して、計算結果を書き足していくのです。
この流れを\term{残差ストリーム}{residual stream}と呼びます。

\begin{figure}[tbp]
\centering
\begin{tikzpicture}[x=9.2mm,y=10mm]
  \draw[line width=5pt,bookblue!35,->] (0,0) -- (11.6,0);
  \node[lbl,anchor=east] at (-0.1,0) {$\vh^{(0)}$};
  \node[lbl,anchor=west] at (11.7,0) {$\vh^{(L)}$};
  \foreach \x/\n/\c/\s in {1.5/注意/redbox/1, 3.8/MLP/greenbox/1, 6.6/注意/redbox/2, 8.9/MLP/greenbox/2} {
    \node[\c,minimum width=13mm,font=\footnotesize] (b) at (\x,1.3) {\n\ $\ell{=}\s$};
    \draw[arr] (\x-0.7,0) to[out=90,in=-90] (b.south west);
    \draw[arr] (b.south east) to[out=-90,in=90] (\x+0.8,0.05);
    \node[op,fill=white,inner sep=0pt,minimum size=4mm,font=\scriptsize] at (\x+0.8,0) {$+$};
  }
  \node[lbl] at (10.5,0.5) {$\cdots$};
\end{tikzpicture}
\caption{残差ストリームの見方。太い青の帯が各位置を流れるベクトル（形は $(B,T,d)$ のまま）で、注意と MLP は帯から読み出し（左の矢印）、計算結果を書き足す（右の矢印と $+$）。
どの部品も帯を「置き換え」ない。損失の勾配も、帯に沿ってそのまま入口まで流れる。}
\label{fig:tf-stream}
\end{figure}

\figref{fig:tf-stream}の太い帯がその流れです。
帯の上の部品は、帯を書き換えずに「メモを書き足す」だけなので、ある部品が役に立たない段階でも、帯の情報は壊れません。
学習の初めに、足し込む量を小さく初期化しておく（\secref{sec:tf-gpt}）のも、この見方からは自然です。
最初は「ほとんど何も書き足さない」ところから始めて、少しずつ役に立つメモを書けるようにするのです。

\subsection*{レイヤー正規化：大きさをそろえる}

残差ストリームには、層を重ねるたびに値が足し込まれていきます。
そのままだと、ベクトルの大きさが層ごとにまちまちになり、注意や MLP は「毎回違う大きさの入力」を相手にしなければなりません。
そこで、部品に入れる前に、ベクトルの大きさをそろえます。
これが\term{レイヤー正規化}{layer normalization}（LayerNorm）\cite{ba2016layer}で、式は\chapref{ch:training}の式\eqref{eq:training-layernorm}と同じです。
1 つの位置のベクトル $\vu\in\R^d$ について、その $d$ 個の成分の平均 $\mu$ と分散 $\sigma^2$ を計算し、
\begin{equation}
\mathrm{LN}(\vu)_j = \gamma_j\,\frac{u_j-\mu}{\sqrt{\sigma^2+\epsilon}} + \beta_j,
\qquad \mu=\frac1d\sum_{j=1}^d u_j,\quad \sigma^2=\frac1d\sum_{j=1}^d (u_j-\mu)^2
\label{eq:tf-layernorm}
\end{equation}
とします。
$\epsilon$ は 0 で割らないための小さな数（$10^{-5}$）、$\vgamma,\vbeta\in\R^d$ は学習されるパラメータ（最初は $\vgamma=\vone$、$\vbeta=\vzero$）です。

小さな例で思い出しておきます。
$\vu=(1,2,3,6)$ の平均は 3、平均との差は $(-2,-1,0,3)$、分散は $(4+1+0+9)/4=3.5$、標準偏差は $\sqrt{3.5}\approx1.871$ です。
差を標準偏差で割ると、$\mathrm{LN}(\vu)\approx(-1.069,\,-0.535,\,0,\,1.604)$ になります。
$\vu$ を 10 倍した $(10,20,30,60)$ を入れても、平均 30 を引いて標準偏差 18.71 で割るので、出力は同じです（\figref{fig:tf-ln-example}）。
入力の「大きさ」が変わっても、部品が受け取る値の大きさは一定に保たれるのです。

\begin{figure}[tbp]
\centering
\begin{tikzpicture}[x=1mm,y=1mm,
  v/.style={font=\footnotesize,minimum height=6mm,inner sep=2pt}]
  \node[box,v,minimum width=29mm] (a) at (0,6) {$(1,\ 2,\ 3,\ 6)$};
  \node[orangebox,v,minimum width=29mm] (b) at (0,-6) {$(10,\ 20,\ 30,\ 60)$};
  \node[graybox,v,minimum width=22mm,minimum height=10mm] (ln) at (41,0) {LayerNorm};
  \node[greenbox,v,minimum width=46mm] (o) at (86,0) {$(-1.069,\ -0.535,\ 0,\ 1.604)$};
  \draw[arr] (a.east) -- (ln.west |- 0,2);
  \draw[arr] (b.east) -- (ln.west |- 0,-2);
  \draw[arr] (ln) -- (o);
  \node[lbl,anchor=south] at (0,10) {大きさの違う入力};
  \node[lbl,anchor=south] at (86,4) {同じ出力（平均 0、分散 1）};
\end{tikzpicture}
\caption{レイヤー正規化は、ベクトルの大きさの違いを消す（$\vgamma=\vone$、$\vbeta=\vzero$ のとき）。10 倍した入力を入れても、出力は変わらない。}
\label{fig:tf-ln-example}
\end{figure}

正規化は\textbf{各位置のベクトルごと}に行い、位置どうし・バッチの中のデータどうしは混ぜません。
平均も分散も、その位置の $d$ 個の成分だけから計算するので、未来のトークンの情報が正規化を通って漏れることもありません。
最後の $\vgamma,\vbeta$ は、「そろえた後に、成分ごとに好きな大きさと位置へ戻してよい」という自由度で、パラメータは $2d$ 個です。

\subsection*{どこに正規化を置くか：Post-LN と Pre-LN}

残差接続とレイヤー正規化を組み合わせる順番には、2 通りあります（\figref{fig:tf-preln}）。

\begin{itemize}
\item \textbf{Post-LN}（原論文\cite{vaswani2017attention}）: 足した\textbf{後}に正規化する。$\vh\leftarrow\mathrm{LN}(\vh+F(\vh))$。
\item \textbf{Pre-LN}（GPT-2\cite{radford2019language} 以降）: 部品に入れる\textbf{前}に正規化し、残差の道には何も挟まない。$\vh\leftarrow\vh+F(\mathrm{LN}(\vh))$。\chapref{ch:training}の式\eqref{eq:training-preln}（事前正規化）と同じ形です。
\end{itemize}

\begin{figure}[tbp]
\centering
\begin{tikzpicture}[node distance=3mm,
  s/.style={box,font=\footnotesize,minimum width=18mm,minimum height=6mm}]
  % Post-LN
  \node[lbl] (t1) at (0,0) {入力 $\vh$};
  \node[s,redbox,above=5mm of t1] (f1) {$F$（注意/MLP）};
  \node[op,above=3mm of f1] (p1) {$+$};
  \node[s,graybox,above=3mm of p1] (n1) {LayerNorm};
  \node[lbl,above=3mm of n1] (o1) {出力};
  \draw[arr] (t1) -- (f1); \draw[arr] (f1) -- (p1); \draw[arr] (p1) -- (n1); \draw[arr] (n1) -- (o1);
  \draw[arr] ($(t1.north)+(0,1.5mm)$) -- ++(14mm,0) |- (p1);
  \node[font=\small\bfseries,above=1mm of o1] {Post-LN（原論文）};
  % Pre-LN
  \begin{scope}[xshift=55mm]
  \node[lbl] (t2) at (0,0) {入力 $\vh$};
  \node[s,graybox,above=5mm of t2] (n2) {LayerNorm};
  \node[s,redbox,above=3mm of n2] (f2) {$F$（注意/MLP）};
  \node[op,above=3mm of f2] (p2) {$+$};
  \node[lbl,above=3mm of p2] (o2) {出力};
  \draw[arr] (t2) -- (n2); \draw[arr] (n2) -- (f2); \draw[arr] (f2) -- (p2); \draw[arr] (p2) -- (o2);
  \draw[arr,bookblue,very thick] ($(t2.north)+(0,1.5mm)$) -- ++(14mm,0) |- (p2);
  \node[font=\small\bfseries,above=1mm of o2] {Pre-LN（GPT-2・本書）};
  \node[lbl,anchor=west,align=left] at ($(p2)+(16mm,-6mm)$) {残差の道に\\何もない};
  \end{scope}
\end{tikzpicture}
\caption{Post-LN と Pre-LN。左の Post-LN では、残差の道の途中に LayerNorm が挟まる。右の Pre-LN では、残差の道（青の太線）は入口から出口まで足し算だけでつながる。}
\label{fig:tf-preln}
\end{figure}

\figref{fig:tf-preln}の右の青い太線を見てください。
Pre-LN では、残差の道に足し算しかないので、式\eqref{eq:tf-residual-grad}の近道が、出口から入口まで途切れずにつながります。
Post-LN では、この道の途中に毎回 LayerNorm が挟まり、「何も掛けられない近道」ではなくなります。
Xiong らは解析と実験で、Post-LN では学習の初めに、出口に近い層の勾配が大きくなることを示しました\cite{xiong2020layer}。
導出は本書の範囲を超えるので、図の違いと次の結論だけ押さえれば十分です（比べる実験は演習問題にあります）。
勾配が大きいと、最初の数回の更新でパラメータが動きすぎて、学習が不安定になります。
それを防ぐために、学習率を小さな値から始めるウォームアップ（\secref{sec:training-schedule}）が欠かせません。
Pre-LN ではこの偏りが小さく、ウォームアップに頼る度合いが下がります。
そのため、GPT-2 以降の LLM のほとんどが Pre-LN を使います。

ただし Pre-LN では、残差ストリームそのものは正規化されないまま出口に届きます。
Pre-LN の LayerNorm が大きさをそろえるのは、注意や MLP に渡す\textbf{写し}だけで、ストリームの本体には $2L$ 回の足し込みがそのまま積もるからです（式\eqref{eq:tf-stream}）。
これをそのまま出力ヘッドに入れると、ロジットが大きくなりすぎて softmax の出力が極端になり、学習が不安定になります。
そこで、最後のブロックの後に\textbf{最後の LayerNorm}（\figref{fig:tf-overview}の上のほう）を置き、出力ヘッドに入れる前に大きさをそろえます。
Post-LN では、各ブロックの出口で毎回そろえているので、この追加は要りません。
\chapref{ch:training}の実験でも、この最後の正規化を抜くと、50 枚の MLP は 5 回とも発散しました（\secref{sec:training-residual}）。

\subsection*{Transformer ブロックのコード}

ここまでの部品を 1 つにまとめたのが、Transformer ブロックです。

\snippet{共有ライブラリの Transformer ブロック（Pre-LN）}{code:tf-block}{llmbook/model.py}{block}

6〜9 行目で、2 つの LayerNorm、注意（\coderef{code:attn-class}）、MLP（\coderef{code:tf-mlp}）を用意します。
\code{nn.LayerNorm(config.n_embd)} は、最後の軸（大きさ $d$）について式\eqref{eq:tf-layernorm}を計算し、$\vgamma$（\code{weight}）と $\vbeta$（\code{bias}）を持ちます。
12 行目が「正規化してから注意に通し、結果を残差に足す」、13 行目が「正規化してから MLP に通し、結果を残差に足す」です。
\figref{fig:tf-preln}右の Pre-LN を、そのまま 2 回書いた形になっています。
入力も出力も $(B,T,d)$ なので、このブロックは何個でも重ねられます。

\subsection*{実験：残差接続と正規化を外すと}
\skimread{押さえる結論は「残差接続を外すと、損失がユニグラムの高さで止まる。正規化を外すと、下がり方が遅くなる」です。\figref{fig:tf-ablation}だけ見れば足ります。}

残差接続とレイヤー正規化は、本当に必要なのでしょうか。
差が見えやすいように基盤モデルの 2 倍の 8 層にした小さなモデル（$d=64$、文字単位、漱石コーパスの先頭 50 万字）で、次の 3 つを同じ条件で 600 ステップ学習させて確かめます（\file{ch10/ablation.py}、3 モデル合わせて約 4 分）。
\begin{itemize}
\item \textbf{そのまま}: \coderef{code:tf-block}のとおり。
\item \textbf{残差なし}: 12〜13 行目の \code{x = x + ...} を \code{x = ...} に変えて、足し込まずに置き換える。
\item \textbf{正規化なし}: LayerNorm を、入力をそのまま返す層 \code{nn.Identity()} に置き換える。
\end{itemize}
「正規化なし」では、ブロックの中の 2 つだけでなく、\textbf{最後の LayerNorm も}外しています（正規化がまったくないモデルと比べるためです）。

どのモデルも、ふつうの \code{GPT} を作ってからブロックや正規化の層を差し替え、\code{GPT} と同じ規則（\secref{sec:tf-gpt}）で初期化し直して、条件をそろえています。
訓練損失の推移が\figref{fig:tf-ablation}です。

\begin{figure}[tbp]
\centering
\begin{tikzpicture}
\begin{axis}[width=0.92\linewidth,height=0.35\linewidth,xlabel={学習ステップ},ylabel={訓練損失（10 ステップ平均）},
  xmin=0,xmax=600,ymin=3.8,ymax=8.5,legend style={at={(0.98,0.96)},anchor=north east,font=\footnotesize}]
\addplot[bookblue,thick,no marks] table[col sep=comma,x=step,y=full]{code/ch10/out/ablation.csv};
\addlegendentry{そのまま（残差＋Pre-LN）}
\addplot[bookorange,thick,no marks] table[col sep=comma,x=step,y=no_layernorm]{code/ch10/out/ablation.csv};
\addlegendentry{正規化なし}
\addplot[bookaccent,thick,no marks] table[col sep=comma,x=step,y=no_residual]{code/ch10/out/ablation.csv};
\addlegendentry{残差なし}
\addplot[bookgray,dotted,thick,domain=0:600] {5.352};
\addlegendentry{ユニグラム 5.352}
\end{axis}
\end{tikzpicture}
\caption{8 層の文字単位モデルで、残差接続と LayerNorm を外したときの学習曲線（\file{ch10/out/ablation.csv}）。点線は、同じデータでのユニグラムの損失。残差を外す（赤）とユニグラムの高さで止まり、正規化を外す（橙）と下がり方が遅くなる。}
\label{fig:tf-ablation}
\end{figure}

最後の 50 ステップの平均損失（\file{ch10/out/ablation.txt}）は、「そのまま」が 4.165、「残差なし」が 5.384、「正規化なし」が 4.721 でした。
図の点線は、同じデータでのユニグラム（文字の出やすさだけで当てるモデル、\chapref{ch:lm}）の損失 5.352 です（\file{ch10/ablation_baseline.py} で測りました）。
「残差なし」は、100 ステップを過ぎたあたりから、この線とほぼ同じ 5.3〜5.4 で横ばいです。
文字の出やすさは覚えても、文脈をまだ使えていません。
8 層を置き換えでつなぐと、勾配が入口近くの層まで届きにくく、文脈を使う学習がなかなか進まないのです。
「正規化なし」は、最初の 30 ステップほどは「そのまま」より速く下がりますが（その理由はここでは追いません）、その後は下がり方が鈍り、600 ステップでは 0.56 の差（パープレキシティで約 1.7 倍）がつきます。

\chapref{ch:training}の 50 枚の MLP では、正規化のない残差接続は発散しましたが、ここでは発散しません。
条件が違うからです。
あちらは He の初期化で、足し込む量が元の値と同じくらい大きく、1 段ごとに倍々に膨らみました。
こちらは重みの標準偏差が 0.02 以下と小さく（\secref{sec:tf-gpt}）、学習前に測ると、残差ストリームの値の大きさは 8 ブロックを通っても 1.000 倍のままでした。
出発点では、ストリームに積もる量がごく小さいので、最後の LayerNorm を外してあっても、ロジットは大きくなりません。
発散はしなくても、学習は遅くなるのです。
これは小さなモデルで 600 ステップだけ学習させた結果なので、差の大きさそのものを一般化はできません。
しかし、残差接続と正規化は、深い Transformer を学習させるための前提だと考えてください。

この節で分かったのは、注意と MLP を「残差に足し込む」形でつなぎ、部品の前で LayerNorm をかける（Pre-LN）ことで、深く重ねても学習が進むようになることです。

\section{組み上げる：GPT クラスを読む}\label{sec:tf-gpt}

\begin{intuition}{ここまでの道のり}
注意に足りなかった三つ（\tabref{tab:tf-three}）への答えがそろいました。
\begin{itemize}
\item 語順を知らない → \textbf{位置埋め込み}を入力に足す（式\eqref{eq:tf-input}）。
\item 情報を集めるだけ → \textbf{MLP} で、各位置で加工する（式\eqref{eq:tf-mlp}）。
\item 深く重ねると不安定 → \textbf{残差接続}と\textbf{レイヤー正規化}（\figref{fig:tf-preln}の右）。
\end{itemize}
ここからは後半です。組み立て、学習させ（\secref{sec:tf-train}）、文章を書かせます（\secref{sec:tf-generate}）。
\end{intuition}

部品はすべてそろいました。
位置埋め込み（\secref{sec:tf-position}）、注意（\chapref{ch:attention}）、MLP（\secref{sec:tf-mlp}）、そしてそれらを残差接続とレイヤー正規化でつないだブロック（\secref{sec:tf-residual}）です。
この節の問いは、「これらを\figref{fig:tf-overview}のとおりに組み上げると、コードではどうなるか」です。
本書の共有ライブラリ \code{llmbook.model.GPT} を、設定・初期化・順伝播の順に読みます。

\subsection*{設定と初期化：部品を並べ、重みを決める}

モデルの大きさを決める数（\textbf{ハイパーパラメータ}、\chapref{ch:ml}）は、\code{GPTConfig} というクラスにまとめてあります（\tabref{tab:tf-config}）。
\code{GPTConfig(n_embd=192)} のように、変えたい項目だけを指定して作ります。
この設定を受け取って部品を並べ、重みの初期値を決めるのが、\code{GPT} クラスの前半（\coderef{code:tf-gpt-init}）です。

\begin{table}[tbp]
\centering
\caption{\code{GPTConfig} の項目と、本書で使う 2 つのモデルの値。文字単位モデルは\secref{sec:tf-train}で学習させる。}
\label{tab:tf-config}
\small
\begin{tabular}{lllrr}
\toprule
項目 & 記号 & 意味 & 基盤モデル & 文字単位モデル \\
\midrule
\code{vocab_size} & $V$ & 語彙の大きさ & 4,096 & 4,078 \\
\code{block_size} & $T_{\max}$ & 文脈長の上限 & 128 & 64 \\
\code{n_layer} & $L$ & ブロックの数 & 4 & 4 \\
\code{n_head} & $H$ & 注意のヘッドの数 & 4 & 4 \\
\code{n_embd} & $d$ & 埋め込みの次元 & 192 & 128 \\
\code{dropout} & — & ドロップアウトの率 & 0.0 & 0.0 \\
\bottomrule
\end{tabular}
\end{table}

\snippet{GPT の部品と初期化}{code:tf-gpt-init}{llmbook/model.py}{gpt-init}

5〜6 行目がトークン埋め込み（$V\times d$ の表）と位置埋め込み（$T_{\max}\times d$ の表）です。
7 行目は埋め込みの直後に置くドロップアウトです（率が 0 なので、本書のモデルでは何もしません）。
8〜9 行目で、ブロックを \code{n_layer} 個作って \code{nn.ModuleList} に入れます。
\code{nn.ModuleList} は、ふつうの Python のリストと違い、中の部品のパラメータを PyTorch に登録します。
これで \code{model.parameters()} がブロックの中の重みも含むようになり、オプティマイザが更新できます。
10 行目が最後の LayerNorm、11 行目が出力ヘッド（$d\to V$ の \code{Linear}）です。

\paragraph{重み共有} 12 行目は、出力ヘッドの重みを、トークン埋め込みの重みと\textbf{同じテンソル}にしています（\code{nn.Linear(d, V)} の重みは「出力 × 入力」の $(V,d)$ という形で、埋め込みの表と同じ形なので、代入できます）。
これを\term{重み共有}{weight tying}\cite{press2017using}と呼びます（\figref{fig:tf-tying}）。
意味を考えてみましょう。
トークン埋め込みの表 $\mE\in\R^{V\times d}$ の $v$ 行目 $\ve_v$ は、トークン $v$ を表すベクトルでした。
出力ヘッドも同じ $\mE$ を使うと、位置 $t$ の最後のベクトル $\vh_t$ からロジットを
\begin{equation}
\vz_t = \mE\,\vh_t,\qquad z_{t,v} = \ve_v^\T\vh_t
\label{eq:tf-tying}
\end{equation}
と作ることになります（$\mE$ は初めから $V\times d$ の形なので、転置なしで掛けます）。
トークン $v$ のロジットは、$\vh_t$ とトークン $v$ の埋め込みとの内積、つまり\textbf{似ている度合い}です（\chapref{ch:attention}で、内積を類似度として使いました）。
「次に来そうなトークンの埋め込みに似たベクトルを作れば、そのトークンの確率が高くなる」という、分かりやすい形になります。
たとえば $\ve_1=(1,0)$、$\ve_2=(0,1)$、$\vh_t=(2,-1)$ なら、$z_{t,1}=1\times2+0\times(-1)=2$、$z_{t,2}=0\times2+1\times(-1)=-1$ で、向きの近い $\ve_1$ のほうが大きくなります。
同じ表を入口と出口で使うので、パラメータも $V\times d$ 個（基盤モデルで 786,432 個）減ります（GPT-2 も同じです）。

\begin{figure}[tbp]
\centering
\begin{tikzpicture}[x=1mm,y=1mm,
  b/.style={font=\scriptsize,minimum height=9mm,inner sep=1.5pt,align=center}]
  \node[tok,minimum width=10mm] (id) at (0,0) {ID $v$};
  \node[purplebox,b,minimum width=21mm] (look) at (21,0) {$\mE$ の $v$ 行目\\を取り出す};
  \node[font=\small] (e) at (39,0) {$\ve_v$};
  \node[box,b,minimum width=22mm] (blk) at (57,0) {位置を足し、\\ブロック $\times L$};
  \node[font=\small] (h) at (75,0) {$\vh_t$};
  \node[purplebox,b,minimum width=21mm] (dot) at (93,0) {$\mE$ の各行\\との内積};
  \node[tok,minimum width=10mm] (z) at (112,0) {$\vz_t$};
  \draw[arr] (id) -- (look); \draw[arr] (look) -- (e); \draw[arr] (e) -- (blk);
  \draw[arr] (blk) -- (h); \draw[arr] (h) -- (dot); \draw[arr] (dot) -- (z);
  \draw[bookpurple,very thick] (look.south) -- ++(0,-5) -| (dot.south);
  \node[lbl,anchor=north] at (57,-10.5) {同じ行列 $\mE$（$V\times d$）を、入口と出口の両方で使う};
  \node[lbl,anchor=south] at (21,5) {入口（埋め込み）};
  \node[lbl,anchor=south] at (93,5) {出口（出力ヘッド）};
\end{tikzpicture}
\caption{重み共有。入口では $\mE$ の行を「引く」、出口では $\mE$ の各行と「内積をとる」。トークン $v$ のロジットは、$\vh_t$ と $\ve_v$ がどれだけ似ているかを表す。}
\label{fig:tf-tying}
\end{figure}

\paragraph{初期化} 14 行目の \code{self.apply(self._init_weights)} は、すべての部品に 22〜24 行目の関数を適用します（21 行目の \code{@staticmethod} は、\code{self} を使わない関数に付ける印です）。
部品が \code{Linear} か \code{Embedding} なら（23 行目の \code{isinstance}）、その重みを平均 0・標準偏差 0.02 の正規分布で初期化します。
0.02 は GPT-2 の実装が使った値です。
\chapref{ch:training}の He の初期化なら $\sqrt{2/192}\approx0.10$ ですが、各部品の前の LayerNorm が入力の大きさをそろえるので、初期値を厳密に合わせなくても学習が始まります。
LayerNorm の $\vgamma$ は 1、$\vbeta$ は 0 のままなので、学習の初めは「平均 0・分散 1 にそろえるだけ」の層です。

16〜19 行目は、\code{named_parameters()} が順に返す（名前, パラメータ）の組のうち、名前が \code{proj.weight} で終わる重み、つまり\textbf{残差ストリームに足し込む直前の層}（注意の出力射影と MLP の 2 層目）だけを、さらに $1/\sqrt{2L}$ 倍小さく初期化し直します（細かい工夫なので、初めて読むときは結論だけで構いません）。
式\eqref{eq:tf-stream}のとおり、残差ストリームには $2L$ 個の量が足し込まれます。
足し込む量がそれぞれ独立で分散 $s^2$ なら、和の分散は $2L\cdot s^2$ に膨らみます（独立な値を足すと分散も足し算になるのでした。\secref{sec:training-init}）。
重みを $1/\sqrt{2L}$ 倍にすれば、その層の出力である足し込む量も $1/\sqrt{2L}$ 倍、分散は $1/(2L)$ 倍になり、和の分散は層の数 $L$ によらず $s^2$ 程度に保たれます。
基盤モデル（$L=4$）では $0.02/\sqrt{8}\approx 0.0071$ です。
これは GPT-2 の論文に書かれた工夫です\cite{radford2019language}。

26〜27 行目の \code{num_params} は、全パラメータの個数を数えます。
\code{self.parameters()} は同じテンソルを 2 回は返さないので、重み共有した $\mE$ は 1 回だけ数えられます。

\subsection*{順伝播}

\snippet{GPT の順伝播}{code:tf-gpt-forward}{llmbook/model.py}{gpt-forward}

6 行目で入力の形 $(B,T)$ を取り出し、7 行目で $T$ が文脈長の上限を超えていないことを確かめます。
位置埋め込みの表は $T_{\max}$ 行しかないので、それより長い入力は扱えないのです。
8 行目の \code{pos} は $(0,1,\dots,T-1)$ という位置の番号の列です（コードの位置は 0 始まり。\code{device=idx.device} は、入力と同じ装置に置く指定で、\chapref{ch:gpu}で使います）。
9 行目が式\eqref{eq:tf-input}で、$(B,T,d)$ のトークン埋め込みに $(T,d)$ の位置埋め込みをブロードキャストで足します。
10〜11 行目でブロックを順に通します。
12 行目で最後の LayerNorm と出力ヘッドを通して、ロジット $(B,T,V)$ を得ます。
14〜16 行目は、正解 \code{targets} が与えられたときに交差エントロピー（\chapref{ch:lm}）を計算します。
\code{targets} は入力を 1 つ先にずらした $(B,T)$ のトークン ID で、どの位置にも「次のトークン」という正解があります。
\code{F.cross_entropy} は、「（個数, クラス数）のロジット」と「（個数,）の正解」の組を受け取るので、バッチの軸と位置の軸を 1 本にまとめます。
15 行目の \code{logits.view(-1, logits.size(-1))} が $(B,T,V)$ を $(BT, V)$ に、16 行目の \code{targets.view(-1)} が $(B,T)$ を $(BT,)$ に並べ直し（たとえば $(4,16,4096)$ を $(64,4096)$ に）、$BT$ 個の位置の次トークン予測の損失を平均します。
\code{targets} を渡さなければ、13 行目の \code{loss = None} のまま、ロジットだけが返ります（生成のときの使い方です）。
全位置の予測を一度に学習できるのは、因果マスクのおかげでした（\chapref{ch:attention}）。
\code{ignore_index=-100} は、正解が $-100$ の位置を損失から外す指定です（\chapref{ch:attention}の実験でも使いました）。

\subsection*{形を確かめる}

\tabref{tab:tf-shapes}の形を、実際のモデルで確かめましょう。
PyTorch には、部品に「呼ばれるたびに実行する関数」を取り付ける\textbf{フック}（hook）という仕組みがあります。
これで各部品の入力と出力の形を記録し、$B=2$, $T=8$ の入力を流しました（\file{ch10/count_params.py}）。
次の出力は、「部品の名前、入力の形 \code{->} 出力の形」を、順伝播で計算が終わった順に並べたものです。

\outputfile{ch10/out/shapes.txt}

1 行目で $(2,8)$ のトークン ID が $(2,8,192)$ になり、2 行目で位置埋め込みが $(8,192)$ を出します。
3〜7 行目は最初のブロック（\code{blocks.0}）の中で、MLP の中間（5 行目の \code{mlp.fc}）だけが $4d=768$ 次元に広がり、ほかはすべて $(2,8,192)$ です。
8 行目が最後の LayerNorm（\code{ln_f}）です。
9 行目で最後に $(2,8,4096)$ のロジットになります。
\tabref{tab:tf-shapes}と一致しました。

この節で分かったのは、\code{GPT} クラスが\figref{fig:tf-overview}をほぼそのままコードにしたもので、加えて「重み共有」と「足し込む層を小さく初期化する」という 2 つの工夫を持つことです。

\section{パラメータを数える}\label{sec:tf-params}
\skimread{押さえる結論は「パラメータ数は式\eqref{eq:tf-params}で数えられる。基盤モデルは 2,583,936 個で、半分近くが MLP にある」です。細部は、\chapref{ch:gpu}でメモリを見積もるときに読めば足ります。}

モデルの「大きさ」は、ふつうパラメータの個数で表します。
この節の問いは、「その数は、設定 $V, T_{\max}, L, d$ からどう決まるのか」です。
式で数えられるようになると、設定を見ただけでモデルの大きさ、必要なメモリ、計算量が見積もれるようになります（\chapref{ch:gpu}、\chapref{ch:infra}）。

\subsection*{部品ごとに数える}

一つずつ数えます。
\texttt{Linear} の重みは「入力の次元 × 出力の次元」個、\texttt{Embedding} は「行数 × $d$」個、LayerNorm は $\vgamma$ と $\vbeta$ で $2d$ 個です。
\code{llmbook} の \code{Linear} はすべてバイアスなしでした。

\begin{itemize}
\item トークン埋め込み: $V\times d$。位置埋め込み: $T_{\max}\times d$。
\item ブロック 1 個:
  \begin{itemize}
  \item 注意の \code{qkv}（$d\to 3d$）: $3d^2$。出力射影 \code{proj}（$d\to d$）: $d^2$。合わせて $4d^2$（\chapref{ch:attention}）。
  \item MLP の \code{fc}（$d\to 4d$）: $4d^2$。\code{proj}（$4d\to d$）: $4d^2$。合わせて $8d^2$。
  \item LayerNorm 2 個: $2\times 2d = 4d$。
  \end{itemize}
  合計 $12d^2+4d$。
\item 最後の LayerNorm: $2d$。出力ヘッド: 重み共有なので $0$。
\end{itemize}

全体は
\begin{equation}
N = Vd + T_{\max}d + L\,(12d^2+4d) + 2d
\label{eq:tf-params}
\end{equation}
です。
基盤モデル（$V=4096$, $T_{\max}=128$, $L=4$, $d=192$）を代入します。
$Vd=4096\times192=786{,}432$、$T_{\max}d=128\times 192=24{,}576$ です。
ブロック 1 個は $12\times192^2+4\times192=442{,}368+768=443{,}136$ で、4 個で $1{,}772{,}544$ です。
最後の LayerNorm が $384$ です。
足すと $786{,}432+24{,}576+1{,}772{,}544+384=2{,}583{,}936$ になります。

\subsection*{コードの値と比べる}

手で数えた値を、実際のモデルと比べます。
\file{ch10/count_params.py} は、式\eqref{eq:tf-params}の各項を計算し、基盤モデルと同じ設定で作った \code{GPT} の \code{num_params()} の値と並べて表示します。

\outputfile{ch10/out/count_params.txt}

7 行目の式による合計と、8 行目の \code{model.num_params()} が、どちらも 2,583,936 で一致しました。
内訳を\tabref{tab:tf-params}にまとめました。

\begin{table}[tbp]
\centering
\caption{基盤モデル（$V=4096$, $T_{\max}=128$, $L=4$, $d=192$）のパラメータ数の内訳。式\eqref{eq:tf-params}の各項。}
\label{tab:tf-params}
\small
\begin{tabular}{llrr}
\toprule
部品 & 式 & 個数 & 割合 \\
\midrule
トークン埋め込み（＝出力ヘッド） & $Vd$ & 786,432 & 30.4\% \\
位置埋め込み & $T_{\max}d$ & 24,576 & 1.0\% \\
注意（4 層） & $L\cdot 4d^2$ & 589,824 & 22.8\% \\
MLP（4 層） & $L\cdot 8d^2$ & 1,179,648 & 45.7\% \\
LayerNorm（9 個） & $L\cdot 4d+2d$ & 3,456 & 0.1\% \\
\midrule
合計 & 式\eqref{eq:tf-params} & 2,583,936 & 100\% \\
\bottomrule
\end{tabular}
\end{table}

\tabref{tab:tf-params}の割合を見てください。
いちばん大きいのは MLP で、注意の 2 倍です（$8d^2$ 対 $4d^2$）。
「Transformer は注意のモデル」という印象とは裏腹に、基盤モデルでは全体の約 46\%、つまり半分近くが MLP にあるのです。
ブロックの部分（$12Ld^2$）だけで見れば、MLP は 3 分の 2 を占めます。
また、本書のモデルは小さいので、埋め込みが 3 割を占めています。
$d$ が大きくなると $d^2$ に比例するブロックの項が支配的になり、埋め込みの割合は小さくなります（\figref{fig:tf-param-bar}）。

\begin{figure}[tbp]
\centering
\begin{tikzpicture}
\begin{axis}[xbar stacked,width=0.84\linewidth,height=34mm,bar width=5.5mm,
  xmin=0,xmax=100,xlabel={パラメータ数に占める割合（\%）},
  ytick={0,1},yticklabels={基盤モデル,GPT-2（最大）},enlarge y limits=0.55,
  ymajorgrids=false,area legend,legend columns=4,
  legend style={at={(0.5,1.06)},anchor=south,font=\footnotesize,/tikz/every even column/.append style={column sep=2mm}}]
\addplot[fill=figpurple,draw=bookpurple] table[col sep=comma,x=emb,y expr=\coordindex]{code/ch10/out/param_share.csv};
\addlegendentry{トークン埋め込み}
\addplot[fill=figorange,draw=bookorange] table[col sep=comma,x=pos,y expr=\coordindex]{code/ch10/out/param_share.csv};
\addlegendentry{位置埋め込み}
\addplot[fill=figred,draw=bookaccent] table[col sep=comma,x=attn,y expr=\coordindex]{code/ch10/out/param_share.csv};
\addlegendentry{注意}
\addplot[fill=figgreen,draw=bookgreen] table[col sep=comma,x=mlp,y expr=\coordindex]{code/ch10/out/param_share.csv};
\addlegendentry{MLP}
\end{axis}
\end{tikzpicture}
\caption{パラメータの内訳の割合（\file{ch10/out/param_share.csv}。式\eqref{eq:tf-params}による。0.1\% 以下の LayerNorm は省いた）。
基盤モデル（258 万個）では埋め込みが 3 割を占めるが、GPT-2 の最大のモデル（約 15.6 億個）では 5\% になり、注意と MLP がほぼすべてを占める。MLP は注意の 2 倍。}
\label{fig:tf-param-bar}
\end{figure}

\begin{quiz}
基盤モデルの $d$ を 192 から 2 倍の 384 にすると、ブロック 1 個のパラメータ数はおよそ何倍になるでしょうか。$L$ を 2 倍にした場合と比べてください。
\end{quiz}

\paragraph{答え}
ブロック 1 個は $12d^2+4d$ なので、$d$ を 2 倍にするとおよそ $2^2=4$ 倍（$12\times384^2+4\times384=1{,}771{,}008$、元の約 4.0 倍）になります。
$L$ を 2 倍にしたときは、ブロック全体がちょうど 2 倍になるだけです。
幅（$d$）を広げるほうが、深さ（$L$）を増やすよりも速くパラメータが増えるのです。

\subsection*{本物の LLM と比べる}

同じ式\eqref{eq:tf-params}で、GPT-2 の最小モデル（$V=50{,}257$, $T_{\max}=1024$, $L=12$, $d=768$）も数えられます。
GPT-2 は \code{Linear} にバイアスがあるので、その分（1 層あたり $3d+d+4d+d=9d$）を足すと、約 1.24 億になります（演習問題）。
\tabref{tab:tf-scale}に、本書のモデルと GPT-2、GPT-3 の設定を並べました。

\begin{table}[tbp]
\centering
\caption{本書のミニGPT と、同じ設計（デコーダのみ・学習型の位置埋め込み・Pre-LN）の実際のモデルの規模。GPT-2 は公開モデルの値~\cite{radford2019language}、GPT-3 は\cite{brown2020language}による。}
\label{tab:tf-scale}
\small
\begin{tabular}{lrrrrrr}
\toprule
モデル & 層 $L$ & 次元 $d$ & ヘッド $H$ & 文脈長 & 語彙 $V$ & パラメータ \\
\midrule
本書の基盤モデル & 4 & 192 & 4 & 128 & 4,096 & 258 万 \\
本章の文字単位モデル & 4 & 128 & 4 & 64 & 4,078 & 132 万 \\
GPT-2（最小） & 12 & 768 & 12 & 1,024 & 50,257 & 1.24 億 \\
GPT-2（最大） & 48 & 1,600 & 25 & 1,024 & 50,257 & 15.6 億 \\
GPT-3 & 96 & 12,288 & 96 & 2,048 & 50,257 & 1,750 億 \\
\bottomrule
\end{tabular}
\end{table}

GPT-3 は本書の基盤モデルの約 7 万倍の大きさですが（$1{,}750\text{ 億}\div258\text{ 万}\approx6.8\text{ 万}$）、部品の種類も組み方も、本章で作ったものと本質的に同じです。
違うのは $L$ と $d$ の値、そして学習に使うデータと計算の量です（\chapref{ch:pretraining}、\chapref{ch:infra}）。

この節で分かったのは、モデルが大きくなるほどパラメータ数は式\eqref{eq:tf-params}の $12Ld^2$ でほぼ決まり、その 3 分の 2（$8Ld^2$）が MLP だということです。

\section{学習させる：文字単位のミニGPT}\label{sec:tf-train}

ここまでで、注意に足りなかった三つ（語順・計算・深さ）を補う部品を作り、それらを組み上げた \code{llmbook.model.GPT} を読み、パラメータを手で数えました。
しかし、組み上げたモデルは、まだ何も学習していません。
この節の問いは、「このモデルを漱石のコーパスで学習させると、損失はどこまで下がり、生成される文章はどう変わっていくか」です。
本格的な事前学習（BPE トークン、30〜40 分）は\chapref{ch:pretraining}で行うので、ここでは CPU で 10 分足らずで終わる\textbf{文字単位}のミニGPT で確かめます。
文字単位にするのは、「1 文字ずつ書く」という分かりやすい設定で、\chapref{ch:lm}のバイグラムやトライグラムと直接比べるためです。

\subsection*{データ、モデル、比べる相手}

データの準備は\chapref{ch:lm}と同じです（\file{ch10/train_char_gpt.py}）。
漱石のコーパスを \code{CharTokenizer}（\chapref{ch:tokenizer}）で 1 文字ずつ番号に変え、末尾 5\%（約 12 万字）を検証用に分けます。
語彙は 4,078 字、訓練データは約 229 万字です。
モデルは\tabref{tab:tf-config}の右の列のとおり、4 層・4 ヘッド・$d=128$、文脈長 64 字です。
基盤モデルより $d$ と文脈長を小さくして、CPU での学習を軽くしました。
パラメータ数は、式\eqref{eq:tf-params}で $4078\times128+64\times128+4\times(12\times128^2+4\times128)+2\times128 = 1{,}318{,}912$ です。
4 つの項は、順にトークン埋め込み $521{,}984$、位置埋め込み $8{,}192$、ブロック 4 個 $788{,}480$、最後の LayerNorm $256$ です（数え方は\secref{sec:tf-params}）。

学習の前に、「どのくらい下がれば上出来か」の目安も用意します。
\chapref{ch:lm}で作った、何も知らないモデル（一様分布。損失は $\log 4078=8.313$）、ユニグラム、バイグラムです。
同じ検証データでの損失は、ユニグラムが 5.198、バイグラムが 3.677 でした（バイグラムは平滑化 $\alpha=0.1$ の値で、スクリプトの中で数え直しています）。

\subsection*{学習ループ}

学習ループは、\chapref{ch:training}で学んだ定石をそのまま使います。

\snippet{文字単位のミニGPT の学習ループ}{code:tf-char-loop}{ch10/train_char_gpt.py}{train-loop}

長く見えますが、骨組みは\chapref{ch:nn}の「5 行のループ」のままです（23 行目が順伝播と損失、24 行目が勾配を消す、25 行目が逆伝播、27 行目が更新）。
ほかの行は、学習を安定させる工夫（\chapref{ch:training}）と途中経過の記録なので、\chapref{ch:training}をまだ読んでいない人は、工夫の中身を気にせず先へ進んで構いません。

1 行目で、1,500 ステップ、バッチ 32、最大学習率 $2\times10^{-3}$ と決めます。
2 行目の \code{make_optimizer} は AdamW（\chapref{ch:training}）を作ります。
5 行目の \code{range(max_steps + 1)} が 1 回多く回るのは、最後の更新のあとにも損失を測るためで、測り終えたら 18〜19 行目の \code{break} で抜けます。
6〜12 行目で、100 ステップごとに訓練損失と検証損失をそれぞれ 20 バッチの平均で見積もります（\code{estimate_loss}）。
13〜17 行目で、決まったステップでは「　吾輩は」の続きを 40 文字生成して残します（温度 0.8、top-$k$ 50。生成の設定は\secref{sec:tf-generate}で説明します）。
14 行目で毎回同じ乱数の種を設定しているのは、どのステップでも同じ「くじの列」を使い、生成例の違いがモデルの違いだけを映すようにするためです。
20〜21 行目で、学習率をウォームアップ 100 ステップ＋コサイン減衰で変えます（\code{lr_at}）。
22〜27 行目が 1 ステップの学習です。
ランダムな位置から 64 文字の窓を 32 本切り出し（\code{get_batch}。正解 \code{y} は、入力 \code{x} を 1 文字先にずらした同じ形の列です）、順伝播で損失を求め、前の勾配を消してから逆伝播し、勾配のノルムを 1.0 で切りそろえて（勾配クリッピング）、パラメータを更新します。

\subsection*{結果}

いよいよ動かします。
\file{code} ディレクトリで \code{python ch10/train_char_gpt.py} と打つと、手元の CPU でも同じ学習が走ります。
著者の環境での実行の記録（\file{ch10/out/train_char_gpt.txt}）から、最初の 3 行と最後の 2 行を抜き出したものが次です。

\begin{tcolorbox}[codebox,colback=outbg,title={実行結果（抜粋）},
  fonttitle=\sffamily\scriptsize\bfseries,coltitle=bookgray,colbacktitle=outbg,
  attach title to upper={\par},top=0mm]
\lstinputlisting[style=booktext,linerange={6-8,21-22}]{code/ch10/out/train_char_gpt.txt}
\end{tcolorbox}

左から、ステップ、訓練損失、検証損失、始めてからの秒数です。
実際には、100 ステップ（30 秒ほど）ごとに 1 行ずつ増えます。
そのたびに損失が小さくなっていく様子を、ぜひ手元の画面で見てください。

数字を読んでいきます。
学習前（step 0）の検証損失は 8.320 で、一様分布の 8.313 とほぼ同じでした。
初期化したばかりのモデルは、どの文字もほぼ同じ確率で予測するからです。
損失はその後順調に下がり、1,500 ステップで訓練 3.129、検証 3.042 に達しました（\figref{fig:tf-curve}）。
所要時間は 474 秒（約 8 分。他の実験と並行して 2 スレッドで実行）でした。

\begin{figure}[tbp]
\centering
\begin{tikzpicture}
\begin{axis}[width=0.92\linewidth,height=0.36\linewidth,xlabel={学習ステップ},ylabel={交差エントロピー},
  xmin=0,xmax=1500,ymin=2.8,ymax=8.6,legend style={at={(0.98,0.96)},anchor=north east,font=\footnotesize}]
\addplot[bookblue,thick,mark=*,mark size=1.2pt] table[col sep=comma,x=step,y=val_loss]{code/ch10/out/train_char_gpt.csv};
\addlegendentry{ミニGPT（検証）}
\addplot[bookgray,thick,dashed,mark=none] table[col sep=comma,x=step,y=train_loss]{code/ch10/out/train_char_gpt.csv};
\addlegendentry{ミニGPT（訓練）}
\addplot[bookaccent,dotted,thick,domain=0:1500] {3.677};
\addlegendentry{バイグラム 3.677}
\addplot[bookorange,dotted,thick,domain=0:1500] {5.198};
\addlegendentry{ユニグラム 5.198}
\end{axis}
\end{tikzpicture}
\caption{文字単位のミニGPT の学習曲線（\file{ch10/out/train_char_gpt.csv}）。点線は数え上げのモデルの検証損失。100 ステップでユニグラムを、400 ステップでバイグラムを下回り、1,500 ステップで 3.04 まで下がる。}
\label{fig:tf-curve}
\end{figure}

\figref{fig:tf-curve}を見ると、最初の 100 ステップで損失が 8.3 から 4.2 へ急に下がっています。
この 100 ステップで、モデルは「よく出る文字」（ユニグラムの情報）を覚え、さらに直前の文字を使った予測も覚え始めて、ユニグラムの 5.198 を下回ります。
その後は下がり方がゆるやかになり、より長い文脈を使った予測を少しずつ身につけていきます。
曲線はまだ下がり続けているので、もっと長く学習させれば、さらに下がるはずです。
なお、検証損失が訓練損失より少し低いのは、検証データ（コーパスの末尾に置かれた作品の一部）が、たまたま訓練データより少し予測しやすかったためで、間違いではありません。

\begin{table}[tbp]
\centering
\caption{同じ検証データ（漱石コーパスの末尾 5\%、約 12 万字）での、1 文字あたりの損失とパープレキシティ（$e^{\text{損失}}$）。
†は 1 トークンあたりの損失（5.512 と 4.863）を、1 トークン = 平均 1.75 字として 1 文字あたりに直した目安の値。}
\label{tab:tf-loss}
\small
\begin{tabular}{lrrl}
\toprule
モデル & 損失 & パープレキシティ & 見ているもの \\
\midrule
一様分布 & 8.313 & 4,078 & 何も見ない \\
ユニグラム & 5.198 & 181 & 文字の頻度 \\
バイグラム & 3.677 & 39.5 & 直前の 1 文字 \\
トライグラム（\chapref{ch:lm}） & 3.398 & 29.9 & 直前の 2 文字 \\
MLP 言語モデル（\chapref{ch:embedding}）† & 約 3.15 & 約 23 & 直前の 3 トークン \\
ミニGPT（300 ステップ） & 3.746 & 42.4 & 直前の最大 64 文字 \\
ミニGPT（1,500 ステップ） & 3.042 & 20.9 & 直前の最大 64 文字 \\
基盤モデル \file{base.pt}† & 約 2.78 & 約 16 & 直前の最大 128 トークン \\
\bottomrule
\end{tabular}
\end{table}

\tabref{tab:tf-loss}は、これまでの章で作ってきた言語モデルと並べたものです。
パープレキシティ（\chapref{ch:lm}）は「平均して何択で迷っているか」でした。
バイグラムは約 40 択、トライグラムは約 30 択で、数えるモデルはここで頭打ちでした（\secref{sec:lm-ngram-limits}）。
8 分学習しただけのミニGPT は、それを下回る約 21 択まで迷いを減らしています。
直前の 1〜2 文字ではなく、64 文字の文脈から必要な情報を注意で集められることが、この差になっています。
\chapref{ch:embedding}の MLP 言語モデルと、30〜40 分かけて学習した基盤モデルは BPE トークン単位なので、1 文字あたりに換算した目安を載せました。
1 トークンの損失は平均 1.75 字ぶんの損失の合計なので、1.75 で割ります（MLP 言語モデルは $5.512/1.75\approx3.15$、基盤モデルは $4.863/1.75\approx2.78$）。
1.75 字は、\chapref{ch:lm}と\chapref{ch:embedding}でも使った、小説の検証データでの値です（\chapref{ch:tokenizer}の 1.68 字は、コーパス全体で測った別の値です）。
単位の違うモデルどうしの比較は大まかなものですが、「見られる文脈が長く、うまく使えるモデルほど損失が低い」という並びは読み取れます。

\subsection*{生成例の変化}

学習の途中で残した生成例が次です。

\outputfile{ch10/out/samples.txt}

step 0 では、ほぼ一様な分布から引くので、コーパスにある 4,078 字がでたらめに並びます。
step 100 では、「あった。」「ない。」「彼の」のように、よく出る仮名と短い言い回しが現れますが、文になっていません。
step 300 では句読点の位置がそれらしくなり、step 1,500 では「その晩は無論」「不思議に対して、こう云う事も」のように、漱石らしい言い回しが続くようになりました。
意味はまだ通りませんが、文字の並びの統計は、確実に学習されています。
読ませたのは 1,500 ステップ × 32 本 × 64 文字で約 307 万文字、コーパスの約 1.3 周分にすぎません。

\paragraph{ここまでで手に入ったもの}
いま動いたのは、借りてきたモデルではありません。
埋め込み（\chapref{ch:embedding}）、注意（\chapref{ch:attention}）、この章の三つの部品を、自分で読んだコードで組んだモデルです。
その約 132 万個のパラメータが、\chapref{ch:nn}から変わらない 5 行のループで 8 分学習しただけで、次のように変わりました。
\begin{itemize}
\item 検証損失は 8.320 から 3.042 へ下がりました（$e^{8.320}\approx4105$ 択の迷いが、$e^{3.042}\approx20.9$ 択に）。
\item \chapref{ch:lm}で頭打ちだった、バイグラム（3.677）とトライグラム（3.398）を下回りました。
\item 書く文は、でたらめな漢字の列から、小説らしい文字の並びになりました。
\end{itemize}
GPT-2 や GPT-3 も、これと同じ種類の部品を、同じ組み方で大きくしたものです（\tabref{tab:tf-scale}）。
ただし、文脈は 64 文字、学習は 8 分だけなので、書けた文の意味はまだ通りません。
もっと大きなモデルを、もっと長く学習させるとどうなるかは、\chapref{ch:pretraining}で確かめます。

この節で分かったのは、組み上げたミニGPT が、CPU で 8 分ほどの学習でバイグラムやトライグラムを下回り、漱石らしい文字の並びを書くようになることです。

\section{文章を生成する：温度・top-\texorpdfstring{$k$}{k}・top-\texorpdfstring{$p$}{p}}\label{sec:tf-generate}

学習したモデルは、文脈を与えると次のトークンの確率分布を出します。
その分布から次のトークンを選ぶ方法は、\secref{sec:lm-generation}で学びました。
貪欲法、温度、top-$k$、top-$p$ です（\tabref{tab:lm-sampling}）。
そこではトライグラムが相手でしたが、選び方は、分布を出すモデルが何であっても同じように使えます。
この節の問いは、「Transformer では、その選び方をどうコードに書くか。設定によって、基盤モデルの書く文章はどう変わるか」です。
ここからは、文字単位モデルではなく、基盤モデル \file{base.pt} を使います。
8 分しか学習していない文字単位モデルでは、設定による文章の違いが見えにくいからです。
\file{base.pt} は、同じ \code{GPT} クラス（\tabref{tab:tf-config}の「基盤モデル」の列）を BPE トークンで 30〜40 分学習させた完成品で、\chapref{ch:attention}でも借りました（学習のやり方は\chapref{ch:pretraining}）。

\subsection*{おさらい：自己回帰と温度}

文章の生成は、(1) いまの文脈をモデルに入れて\textbf{最後の位置}のロジット $\vz\in\R^V$ を得る、(2) $\vz$ から次のトークンを 1 つ選ぶ、(3) 選んだトークンを文脈の末尾に付け足して (1) に戻る、の繰り返し（自己回帰）でした。
(2) で、毎回いちばん確率の高いトークンを選ぶのが貪欲法、確率に従ってくじを引くのがサンプリングです。
貪欲法は、文章では同じ言い回しの繰り返しに陥りやすいことが知られています\cite{holtzman2020curious}。
サンプリングの際、ロジットを温度 $\tau$ で割ってから softmax をとると（式\eqref{eq:lm-temperature}）、分布のとがり具合を調整できるのでした。
$\tau<1$ なら確率の高いトークンに集中し、$\tau>1$ なら平らになります。

基盤モデルに「どうも仕方が」を与えて、確かめましょう（\file{ch10/sampling_demo.py}）。
最後の位置のロジットを 3 通りの温度で割ったときの、上位 8 個の確率が\figref{fig:tf-temp}です（数値は \file{ch10/out/temperature.txt}）。

\begin{figure}[tbp]
\centering
\begin{tikzpicture}
\begin{axis}[ybar,bar width=2.6mm,width=0.92\linewidth,height=0.33\linewidth,
  xlabel={温度 1 での順位},ylabel={確率},xtick={1,...,8},ymin=0,ymax=1,
  legend style={at={(0.98,0.95)},anchor=north east,font=\footnotesize}]
\addplot[fill=bookblue,draw=bookblue] table[col sep=comma,x=rank,y=t05]{code/ch10/out/temperature.csv};
\addlegendentry{$\tau=0.5$}
\addplot[fill=bookorange,draw=bookorange] table[col sep=comma,x=rank,y=t10]{code/ch10/out/temperature.csv};
\addlegendentry{$\tau=1.0$}
\addplot[fill=bookgreen,draw=bookgreen] table[col sep=comma,x=rank,y=t20]{code/ch10/out/temperature.csv};
\addlegendentry{$\tau=2.0$}
\end{axis}
\end{tikzpicture}
\caption{基盤モデルに「どうも仕方が」を与えたときの、次トークンの上位 8 個の確率と温度。1 位は「ない。」。
$\tau=0.5$（青）では 1 位に 9 割以上が集まり、$\tau=2$（緑）では 1 位でも 5.5\% まで下がり、2 位以下との差が小さくなって平らになる。}
\label{fig:tf-temp}
\end{figure}

温度 1.0 での 1 位「ない。」の確率 0.559 は、\chapref{ch:intro}で見た基盤モデルの予測と同じです。
温度 0.5 にすると 0.934 に跳ね上がり、温度 2.0 にすると 0.055 まで下がります。
温度 2.0 では、上位 8 個を合わせても 2 割に届きません（\figref{fig:tf-temp}の緑。8 個の確率を足すと 0.151 です）。
残りの 8 割以上（約 85\%）は、4,000 個以上ある「ありそうもないトークン」に薄く広がっているのです。

\subsection*{おさらい：候補を絞る top-\texorpdfstring{$k$}{k} と top-\texorpdfstring{$p$}{p}}

温度を上げると、そうした「ありそうもないトークン」を引いてしまう危険も増えます。
そこで、くじを引く前に候補を絞ります。
top-$k$ は、確率の高い順に $k$ 個だけを残します。
top-$p$\cite{holtzman2020curious} は、確率の高い順に足していき、合計が初めて $p$ を超えるところまでを残します。
top-$k$ は候補の数を固定しますが、top-$p$ は分布の形に合わせて候補の数が変わります。

top-$p$ をコードに書くときは、「\textbf{自分より上位の確率の和}が $p$ を超えていたら捨てる」と判定します。
自分がいなくても、上位だけで合計がもう $p$ を超えている、ということだからです。
\tabref{tab:tf-topp}の例では、4 番目だけが捨てられ、上位 3 個（合計 0.95）が残ります。

\begin{table}[tbp]
\centering
\caption{top-$p$ の判定の例（$p=0.9$）。確率の高い順に並べ、「自分より上位の和」が $p$ を超えたトークンを捨てる。}
\label{tab:tf-topp}
\small
\begin{tabular}{lrrrr}
\toprule
順位 & 1 & 2 & 3 & 4 \\
\midrule
確率 & 0.5 & 0.3 & 0.15 & 0.05 \\
自分より上位の和 & 0 & 0.5 & 0.8 & 0.95 \\
0.9 を超えるか & いいえ & いいえ & いいえ & はい \\
判定 & 残す & 残す & 残す & 捨てる \\
\bottomrule
\end{tabular}
\end{table}

\begin{quiz}
並べた確率が $(0.4,\,0.3,\,0.2,\,0.1)$ のとき、top-$p$（$p=0.75$）で残るトークンは何個でしょうか。
\tabref{tab:tf-topp}と同じ表を作って確かめてください。
\end{quiz}

\paragraph{答え}
「自分より上位の和」は $(0,\,0.4,\,0.7,\,0.9)$ です。
0.75 を超えるのは 4 番目だけなので、3 個が残ります。
上位 2 個の合計は 0.7 で、まだ $p=0.75$ に届いていないので、3 番目は捨てられません。

\subsection*{generate を読む}
\skimread{押さえる結論は「\code{generate} は、\figref{fig:tf-genflow}の流れを 1 トークンごとにくり返す関数で、温度・top-$k$・top-$p$ は引数で渡す」です。行ごとの読解は 2 周目で足ります。}

\code{llmbook} の \code{generate} は、ここまでの方法をすべて実装しています。
1 トークンを選ぶまでの流れは\figref{fig:tf-genflow}のとおりです。

\begin{figure}[tbp]
\centering
\begin{tikzpicture}[x=1mm,y=1mm,
  b/.style={font=\scriptsize,minimum width=17.5mm,minimum height=9mm,inner sep=1pt,align=center}]
  \node[tok,b] (c) at (0,0) {文脈\\$(B,T)$};
  \node[box,b] (m) at (21.5,0) {モデル\\最後の位置};
  \node[orangebox,b] (t) at (43,0) {$\tau$ で割る};
  \node[greenbox,b] (k) at (64.5,0) {top-$k$\\で絞る};
  \node[greenbox,b] (p) at (86,0) {top-$p$\\で絞る};
  \node[purplebox,b] (s) at (107.5,0) {softmax\\くじを引く};
  \draw[arr] (c) -- (m); \draw[arr] (m) -- (t); \draw[arr] (t) -- (k); \draw[arr] (k) -- (p); \draw[arr] (p) -- (s);
  \draw[arr,bookblue,very thick] (s.south) -- ++(0,-9) -| (c.south);
  \node[lbl,anchor=north] at (53.75,-14.5) {選んだトークンを文脈の末尾に付け足す（31 行目）};
  \node[lbl,anchor=south] at (0,5) {7 行目};
  \node[lbl,anchor=south] at (21.5,5) {8〜9 行目};
  \node[lbl,anchor=south] at (43,5) {16 行目};
  \node[lbl,anchor=south] at (64.5,5) {17〜20 行目};
  \node[lbl,anchor=south] at (86,5) {21〜28 行目};
  \node[lbl,anchor=south] at (107.5,5) {29〜30 行目};
\end{tikzpicture}
\caption{\code{generate} が 1 トークンを選ぶ流れ。箱の上は\coderef{code:tf-generate}の行番号。これを \code{max_new_tokens} 回くり返す。
温度が 0 のときは、「$\tau$ で割る」から先の代わりに、最大のロジットのトークンを選ぶ（10〜15 行目）。}
\label{fig:tf-genflow}
\end{figure}

\snippet{共有ライブラリの生成関数}{code:tf-generate}{llmbook/model.py}{generate}

1 行目の \code{@torch.no_grad()} は、勾配の計算に必要な記録をとらない指定で、推論を速く、省メモリにします。
6 行目のループが、上の (1)〜(3) の繰り返しです。
7 行目で文脈を後ろから \code{block_size} 個に切り詰めます。位置埋め込みの表は 128 行しかないからです。
8〜9 行目でモデルに通し、最後の位置のロジット $(B,V)$ だけを取り出します。
10〜15 行目は温度 0 の場合（貪欲法）で、\code{argmax} で最大のトークンを選んで \code{torch.cat} で末尾に付け足し、終わりの印（\code{stop_token}）がバッチの全部で出たら（\code{.all()}）止めます。
温度 0 だけを別に扱うのは、16 行目の「温度で割る」計算が、0 ではできないからです。
16 行目が、ロジットを温度で割る計算です。
17〜20 行目が top-$k$ で、\code{torch.topk} で $k$ 番目に大きいロジットを求め、それより小さいロジットを $-\infty$（確率 0）にします。

21〜28 行目が top-$p$ です（考え方は\coderef{code:lm-sample-next}と同じなので、細かい書き方は読み飛ばしても構いません）。
22 行目でロジットを大きい順に並べ、23 行目で確率にします。
25 行目の \code{probs.cumsum(dim=-1) - probs} が、\tabref{tab:tf-topp}の「自分より上位の和」です。
これが $p$ を超えるトークンのロジットを $-\infty$ にします。
27〜28 行目の \code{scatter} は、並べ替えたロジットを元の語彙の順に戻す操作です。
29〜31 行目で確率に直し、\code{torch.multinomial} で確率に従って 1 つ引いて、末尾に付け足します。

\subsection*{設定による出力の違い}

基盤モデルに「\texttt{\textbackslash n}　吾輩は」を与え、6 通りの設定で 30 トークンずつ生成させました（\file{ch10/sampling_demo.py} の後半）。
違いが設定だけから生まれるように、どの設定でも乱数の種を同じにしてあります。
改行は「／」で、文字にならないバイト列は「□」で示しています。

\outputfile{ch10/out/sampling.txt}

1 行目の貪欲法では、最初はもっともらしい文が続きますが、最後に「君に君に君に」と同じ言い回しの繰り返しに陥っています。
2 行目の温度 0.5 は、文としてのまとまりがよい一方、「例のごとく」が 2 回出るなど、どこかで見たような言い回しが多くなります。
3 行目の温度 1.0 は、モデルの分布そのままから引いた結果ですが、「開布衍思議」のような意味の通らない字の並びと「□」2 個が混ざっています。
4,096 個の語彙には、1 個ずつの確率は小さくても、合わせると無視できない「ありそうもないトークン」の裾があり、30 トークンも引くと、どこかでその裾を引いてしまうのです。
4 行目の温度 1.5 は裾をさらに厚くするので、「□」は 3 個に増えています。
本書の BPE はバイト単位なので（\chapref{ch:tokenizer}）、漢字 1 字の 3 バイトの途中までしか表さないトークンがあります。
裾から引いたそうしたトークンがちぐはぐに並ぶと、文字にならない並びができてしまいます。
5 行目の top-$k$ 10 と 6 行目の top-$p$ 0.9 は、温度 1.0 のままでも、変な字がほとんどありません。
この裾を、引く前に候補から外したからです。

書き出しにも注目してください。
2 行目と 5 行目は「例のごとく笑」まで、3 行目と 4 行目は「下宿の門を□交る□ユ」まで同じで、6 行目も「下宿の門を」で始まります。
同じ乱数の列を使ったので、分布が似ているあいだは同じトークンが選ばれるのです。

\begin{practice}{どの設定を使うか}
決まった正解がある課題（計算、コード、事実の質問）には、貪欲法か低い温度（0〜0.3 程度）を使います。
本書でも、計算ドリルの正答率は貪欲法で測ります（\code{llmbook.arith.drill_accuracy}）。
文章の創作や、多様な答えを集めたい場面（強化学習で複数の答えを試させる場面、\chapref{ch:rl-llm}）では、温度 0.7〜1.0 と top-$p$ 0.9〜0.95 の組み合わせがよく使われます。
これらは経験的な目安で、モデルや課題によって最適な値は変わります。
\end{practice}

この節で分かったのは、同じモデルでも、確率分布からの選び方（貪欲法・温度・top-$k$・top-$p$）によって、生成される文章の性質が大きく変わることです。

\section{Transformer の系譜：なぜデコーダのみか}\label{sec:tf-family}
\laterread{飛ばしても、この先の章は読めます。「エンコーダ」や BERT という言葉に出会ったときに戻れば足ります。}

作るものはすべてそろったので、最後に、本章の「デコーダのみ」の Transformer の立ち位置を確かめます。
デコーダではないもの、つまり「エンコーダ」とは何でしょうか。
この節の問いは、「原論文の Transformer と本書のミニGPT はどう違い、なぜ現在の LLM はデコーダのみを使うのか」です。

\subsection*{原論文はエンコーダ–デコーダ}

Transformer の原論文\cite{vaswani2017attention}は、英語の文を読んでドイツ語の文を書く、機械翻訳のためのモデルでした。
そのため、二つの部分からできていました（\figref{fig:tf-family}の左）。
翻訳元の文全体を読む\textbf{エンコーダ}と、翻訳先の文を 1 トークンずつ書く\textbf{デコーダ}です。
エンコーダは因果マスクを使わないので、各トークンが文の前後をすべて見られます。
デコーダは、因果マスク付きの自己注意に加えて、エンコーダの出力を見にいく\term{交差注意}{cross-attention}を持ちます。
交差注意は、クエリをデコーダ側のベクトルから、キーとバリューをエンコーダの出力から作る注意で、式は\chapref{ch:attention}の縮小付き内積注意と同じです。

\begin{figure}[tbp]
\centering
\begin{tikzpicture}[s/.style={box,font=\footnotesize,minimum width=19mm,minimum height=8mm,align=center}]
  \node[s,greenbox] (enc) at (0,0) {エンコーダ\\（マスクなし）};
  \node[s,redbox] (dec) at (3.4,0) {デコーダ\\（因果マスク）};
  \draw[arr] (enc) -- node[lbl,above]{交差注意} (dec);
  \node[lbl] at (0,-0.75) {翻訳元を読む};
  \node[lbl] at (3.4,-0.75) {翻訳先を書く};
  \node[s,greenbox] (bert) at (6.8,0) {エンコーダ\\のみ};
  \node[lbl] at (6.8,-0.75) {文全体を読む};
  \node[s,redbox] (gpt) at (9.8,0) {デコーダ\\のみ};
  \node[lbl] at (9.8,-0.75) {続きを書く};
  \foreach \x/\t/\u in {1.7/エンコーダ–デコーダ/原論文 (2017)、T5, 6.8/エンコーダのみ/BERT (2018), 9.8/デコーダのみ/GPT (2018)〜現在の LLM} {
    \node[font=\footnotesize\bfseries] at (\x,1.0) {\t};
    \node[lbl] at (\x,1.45) {\u};
  }
\end{tikzpicture}
\caption{Transformer の 3 つの系統。緑はマスクなしの自己注意（前後を見る）、赤は因果マスク付きの自己注意（過去だけを見る）。本書のミニGPT と現在の LLM の多くは右の「デコーダのみ」。}
\label{fig:tf-family}
\end{figure}

原論文の翌年、その片方だけを使うモデルが相次いで現れました。
エンコーダだけを使い、隠したトークンを前後から当てる「穴埋め」で学習する BERT\cite{devlin2019bert}と、デコーダから交差注意を取り除いて使う GPT\cite{radford2018improving}（後継が GPT-2\cite{radford2019language}）です。
エンコーダ–デコーダの形も、T5\cite{raffel2020exploring}のように、あらゆる課題を「文から文への変換」とみなすモデルとして発展しました。
三つの系統の違いを\tabref{tab:tf-family3}にまとめます。

\begin{table}[tbp]
\centering
\caption{Transformer の 3 つの系統の比較。本書のミニGPT は右端の列。}
\label{tab:tf-family3}
\small
\begin{tabular}{llll}
\toprule
 & エンコーダ–デコーダ & エンコーダのみ & デコーダのみ \\
\midrule
自己注意 & 読む側は前後、書く側は過去 & 前後を見る & 過去だけを見る \\
学習の課題 & 文から文への変換 & 穴埋め & 次トークン予測 \\
向いている用途 & 翻訳、要約 & 文の分類、検索 & 文章の生成 \\
代表 & 原論文、T5 & BERT & GPT、現在の LLM \\
\bottomrule
\end{tabular}
\end{table}

\subsection*{なぜ LLM はデコーダのみが主流か}

現在の LLM の多くは、デコーダのみです。
理由は、厳密に証明されたものではありませんが、次のように説明されます。

第一に、\textbf{学習の目的が単純で、データがいくらでもある}ことです。
次トークン予測は、どんな文章もそのまま教材になります（自己教師あり学習、\chapref{ch:lm}）。
しかも因果マスクのおかげで、1 つの系列の全位置が同時に学習の問題になります。
BERT の穴埋めでは、損失を計算するのは隠した一部（原論文では 15\%）のトークンだけです。

第二に、\textbf{生成と理解が 1 つの形で済む}ことです。
質問への回答も、翻訳も、要約も、「指示の文の続きを書く」という 1 つの形にできます。
GPT-3\cite{brown2020language}は、大きなデコーダのみのモデルが、例をいくつか文脈に並べるだけで新しい課題をこなせることを示しました。

第三に、\textbf{構造が単純で、大きくしやすい}ことです。
部品の種類が少ないので、層と次元を増やすだけで規模を伸ばせます（\chapref{ch:pretraining}、\chapref{ch:infra}）。

ただし、エンコーダのみのモデルも、検索用の文の埋め込みなどで今も広く使われています。
デコーダのみが常に優れているのではなく、大規模な汎用の言語モデルで標準になった、ということです。

この節で分かったのは、本書のミニGPT が、原論文のデコーダから交差注意を取り除き、次トークン予測だけで学習する系統（GPT）に属することです。

\begin{summary}
\item Transformer（デコーダのみ）は、埋め込み＋位置 → ブロック × $L$ → 最後の LayerNorm → 出力ヘッドで、テンソルの形は入口と出口以外 $(B,T,d)$ のまま。
\item 注意は語順を知らないので、位置ごとのベクトルを埋め込みに足す（位置エンコーディング）。本書と GPT-2 は学習型、原論文は正弦波。
\item 注意は位置の間の「通信」、MLP は各位置での「計算」。MLP は $d\to4d\to d$ の 2 層で、ブロックのパラメータの 3 分の 2（基盤モデルでは全体の約 46\%）を占める。
\item 残差接続は、勾配が何も掛けられずに戻る「近道」を作り、深く重ねても学習を可能にする。各位置を流れるベクトルを残差ストリームとみなせる。
\item LayerNorm は各位置のベクトルを平均 0・分散 1 にそろえる。部品の前に置く Pre-LN が GPT-2 以降の標準。
\item 出力ヘッドはトークン埋め込みと重みを共有し、残差に足し込む層は $1/\sqrt{2L}$ 倍小さく初期化する。
\item パラメータ数は $Vd+T_{\max}d+L(12d^2+4d)+2d$。基盤モデルは 2,583,936。
\item 生成では、貪欲法・温度・top-$k$・top-$p$ で選び方を変えられ、文章の性質が大きく変わる。
\end{summary}

\section*{次の章へ}

この章で、自分のミニGPT が動きました。
しかし、書いた文の意味はまだ通りません。
もっと大きなモデルを、BPE トークンで、もっと長く学習させるとどこまで行けるかは、\chapref{ch:pretraining}（事前学習）で確かめます。

本章で組み立てたミニGPT は、2019 年の GPT-2 と同じ設計です。
しかし、それから現在までに、LLM の設計にはいくつもの改良が加えられてきました。
学習型の位置埋め込みは文脈長の上限を超えられず、生成のたびに過去の全トークンを計算し直すのは無駄です（基盤モデルは 129 番目のトークンを入力できず、\coderef{code:tf-generate}の 8 行目は毎回、文脈の全体を通し直していました）。
次の\chapref{ch:modern-arch}では、現在の LLM が GPT-2 とどこが違い、なぜそう変わったのかを、一つずつ実装しながら確かめます。
初めて読む人は、\chapref{ch:modern-arch}と\chapref{ch:gpu}を後に回して、この章の \code{GPT} クラスを BPE トークンで学習させる\chapref{ch:pretraining}へ進んで構いません。

\section*{さらに学ぶために}

\begin{itemize}
\item Vaswani ら「Attention Is All You Need」\cite{vaswani2017attention}：Transformer の原論文。エンコーダ–デコーダ、正弦波の位置エンコーディング、Post-LN。
\item Radford ら「Language Models are Unsupervised Multitask Learners」\cite{radford2019language}：GPT-2。Pre-LN、最後の LayerNorm、残差の初期化の工夫。本書のミニGPT の設計の元。
\item Karpathy「nanoGPT」\cite{karpathy2022nanogpt}と講義「Let's build GPT」\cite{karpathy2022zerotohero}：\code{llmbook.model} とほぼ同じ構造の実装と、それを一から書く動画。同じ範囲を扱う書籍に Raschka\cite{raschka2024build}がある。
\end{itemize}

\begin{exercises}
\exfirst{1 と 4}
\item （手計算）$\vu=(2,4,4,6)$ に LayerNorm（$\vgamma=\vone$, $\vbeta=\vzero$, $\epsilon=0$）を適用した結果を求めてください。また、$\vu$ のすべての成分に 100 を足したベクトルでは結果がどうなるでしょうか。
\item （手計算）式\eqref{eq:tf-params}に GPT-2 最小モデルの設定（$V=50{,}257$, $T_{\max}=1024$, $L=12$, $d=768$）を代入し、\code{Linear} のバイアス（1 層あたり $9d$）も足して、約 1.24 億になることを確かめてください。$12Ld^2$ は全体の何割でしょうか。
\item （コード）\file{ch10/ablation.py} に「Post-LN」の版（\code{x = ln1(x + attn(x))} のように足した後に正規化する）を加えて、学習曲線を比べてください。
\item （コード）\file{ch10/sampling_demo.py} で、計算ドリル「\texttt{\textbackslash n12+35=}」の続きを温度 0、0.8、1.5 で 20 回ずつ生成し、「47」が出た回数を数えて、結果を\figref{fig:tf-temp}の考え方で説明してください。
\item （考察）エンコーダのみのモデル（BERT）で文章を 1 トークンずつ生成しようとすると、何が困るでしょうか。因果マスクの有無と、学習の目的の違いから説明してください。
\end{exercises}
